%============================================================
%  Bd. 38 - Monoide
%============================================================

\documentclass[main.tex]{subfiles}

\ifSubfilesClassLoaded{
  \BandSetupStandalone{B38}
}{
}

\title{Bd. 38 - Monoide}
\author{Martin Kunze}
\date{}
\setcounter{file}{38}

\hypersetup{
  pdftitle={Bd. 38 - Monoide},
  pdfauthor={Martin Kunze}
}

\begin{document}

\title{Bd. 38 - Monoide}
\author{Martin Kunze}
\date{}
\hypersetup{pageanchor=false}
\maketitle
\hypersetup{pageanchor=true}
\setcounter{tocdepth}{3}
\tableofcontents
\setcounter{file}{38}
\renewcommand{\FormulaBandID}{38}

\chapter{Einführung}

Dieser Band entwickelt Monoide als Halbgruppen mit einem neutralen Element.
Darauf aufbauend werden Einheiten, Teilbarkeit, Irreduzibilität,
Primelemente und Morphismen von Monoiden untersucht.

\chapter{Monoide}

\section{Definition des Monoids}

\FormulaDefDeltaK[Begriff des Monoids]{%
  \MonoidAlg{A}{\star}{e}%
}{MonoidDef}{%
  \DeltaRow{Mengen}{A\dsep e\dsep x}
  \DeltaRow{\textbf{Axiome}}{}
  \DeltaPrem{Halbgruppe}{\SemigroupAlg{A}{\star}}
  \DeltaRow{Neutrales Element im Träger}{e\in A}
  \DeltaRow{Linksneutralität}{x\in A\vdash e\star x=x}
  \DeltaRow{Rechtsneutralität}{x\in A\vdash x\star e=x}
  \DeltaRow{\textbf{Neues Symbol}}{}
  \DeltaPrem{Monoid}{\MonoidAlg{A}{\star}{e}}
}

\FormulaAxiomDeltaK[Zugriff auf die Halbgruppe eines Monoids]{%
  \MonoidAlg{A}{\star}{e}\vdash\SemigroupAlg{A}{\star}%
}{MonoidBaseSemigroupAxiom}{%
  \DeltaRow{Mengen}{A\dsep e}
}

\FormulaAxiomDeltaK[Zugriff auf das neutrale Element]{%
  \MonoidAlg{A}{\star}{e}\vdash e\in A%
}{MonoidIdentityCarrierAxiom}{%
  \DeltaRow{Mengen}{A\dsep e}
}

\FormulaAxiomDeltaK[Zugriff auf die Linksneutralität]{%
  \MonoidAlg{A}{\star}{e}\dsep x\in A
  \vdash e\star x=x%
}{MonoidLeftIdentityAxiom}{%
  \DeltaRow{Mengen}{A\dsep e\dsep x}
}

\FormulaAxiomDeltaK[Zugriff auf die Rechtsneutralität]{%
  \MonoidAlg{A}{\star}{e}\dsep x\in A
  \vdash x\star e=x%
}{MonoidRightIdentityAxiom}{%
  \DeltaRow{Mengen}{A\dsep e\dsep x}
}

\section{Das neutrale Element}

\FormulaThmDeltaK[Ein neutraler Kandidat stimmt mit dem neutralen Element überein]{%
  \begin{aligned}[t]
    &\MonoidAlg{A}{\star}{e}\dsep e'\in A\dsep
      \forall x\in A\,(e'\star x=x)\dsep
      \forall x\in A\,(x\star e'=x)\vdash{}\\[-2pt]
    &e=e'
  \end{aligned}
}{MonoidIdentityCandidateEqualsIdentity}{%
  \DeltaRow{Mengen}{A\dsep e\dsep e'\dsep x}
}
\begin{tabproofwide}
    \proofstepwidestar[1]{\MonoidAlg{A}{\star}{e}}{\rA}
    \proofstepwidestar[2]{e'\in A}{\rA}
    \proofstepwidestar[3]{\forall x\in A\,(e'\star x=x)}{\rA}
    \proofstepwidestar[4]{\forall x\in A\,(x\star e'=x)}{\rA}
    \proofstepwidestar[1]{e\in A}{%
      \FormulaRefAuto{MonoidIdentityCarrierAxiom}[ax]{1}}
    \proofstepwidestar[1,3]{e'\star e=e}{\rRE{5,3}}
    \proofstepwide[1,3]{e}{=}{e'\star e}{%
      \FormulaRefAuto{a=b\vdash b=a}{6}}
    \proofstepwide[1,2]{}{=}{e'}{%
      \FormulaRefAuto{MonoidRightIdentityAxiom}[ax]{1,2}}
    \proofstepwidestar[1,2,3]{e=e'}{\rChain{7,8}}
\end{tabproofwide}

\FormulaThmDeltaK[Existenz und Eindeutigkeit des neutralen Elements]{%
  \begin{aligned}[t]
    &\MonoidAlg{A}{\star}{e}\vdash{}\\[-2pt]
    &\exists! e'\in A\,\Bigl(
      \forall x\in A\,(e'\star x=x)\land
      \forall x\in A\,(x\star e'=x)
      \Bigr)
  \end{aligned}
}{MonoidIdentityUnique}{%
  \DeltaRow{Mengen}{A\dsep e\dsep e'\dsep x}
}
\begin{tabproofsplitwide}
  \proofpartwide{Existenz eines neutralen Elements}
    \proofstepwidestar[1]{\MonoidAlg{A}{\star}{e}}{\rA}
    \proofstepwidestar[1]{e\in A}{%
      \FormulaRefAuto{MonoidIdentityCarrierAxiom}[ax]{1}}
    \proofstepwidestar[3]{x\in A}{\rA}
    \proofstepwidestar[1,3]{e\star x=x}{%
      \FormulaRefAuto{MonoidLeftIdentityAxiom}[ax]{1,3}}
    \proofstepwidestar[1,3]{x\star e=x}{%
      \FormulaRefAuto{MonoidRightIdentityAxiom}[ax]{1,3}}
    \proofstepwidestar[1]{\forall x\in A\,(e\star x=x)}{%
      \rUI{\rRI{3,4}}}
    \proofstepwidestar[1]{\forall x\in A\,(x\star e=x)}{%
      \rUI{\rRI{3,5}}}
    \proofstepwidestar[1]{%
      \begin{aligned}[t]
        &e\in A\land\Bigl(
          \forall x\in A\,(e\star x=x)\land{}\\[-2pt]
        &\qquad \forall x\in A\,(x\star e=x)
          \Bigr)
      \end{aligned}}{%
      \rAI{2,\rAI{6,7}}}
  \closeproofpartwide

  \proofpartwide{Eindeutigkeit}
  \setcounter{proofstepnr}{8}
    \proofstepwidestar[9]{%
      \begin{aligned}[t]
        &e'\in A\land\Bigl(
          \forall x\in A\,(e'\star x=x)\land{}\\[-2pt]
        &\qquad \forall x\in A\,(x\star e'=x)
          \Bigr)
      \end{aligned}}{\rA}
    \proofstepwidestar[9]{e'\in A}{\rAEa{9}}
    \proofstepwidestar[9]{\forall x\in A\,(e'\star x=x)}{%
      \rAEa{\rAEb{9}}}
    \proofstepwidestar[9]{\forall x\in A\,(x\star e'=x)}{%
      \rAEb{\rAEb{9}}}
    \proofstepwidestar[1,9]{e=e'}{%
      \FormulaRefAuto{MonoidIdentityCandidateEqualsIdentity}[thm]{%
        1,10,11,12}}
    \proofstepwidestar[1]{%
      \begin{aligned}[t]
        &\exists! e'\in A\,\Bigl(
          \forall x\in A\,(e'\star x=x)\land{}\\[-2pt]
        &\qquad \forall x\in A\,(x\star e'=x)
          \Bigr)
      \end{aligned}}{\UEI{8,9,13}}
  \closeproofpartwide
\end{tabproofsplitwide}

\FormulaThmDeltaK[Potenzhalbgruppe eines Monoids]{%
  \MonoidAlg{A}{\star}{e}\vdash
  \MonoidAlg{\PowNE{A}}{\star_{\mathcal P}}{\{e\}}%
}{PowerSemigroupMonoid}{%
  \DeltaRow{Mengen}{A\dsep e\dsep X}
  \DeltaRow{Funktionen}{\star\dsep\star_{\mathcal P}}
}
\begin{proof}
Die Potenzhalbgruppe ist nach
\FormulaRefAuto{NonemptyPowerSemigroup}[thm] eine Halbgruppe.  Aus
\(e\in A\) folgt \(\{e\}\in\PowNE A\), und für jedes
\(X\in\PowNE A\) gilt
\[
  \{e\}\star_{\mathcal P}X
  =\{e\star x\mid x\in X\}=X,
  \qquad
  X\star_{\mathcal P}\{e\}
  =\{x\star e\mid x\in X\}=X.
\]
Damit ist \(\{e\}\) das neutrale Element der Potenzhalbgruppe.
\end{proof}

\FormulaThmDeltaK[Neutrale Elemente von Potenzhalbgruppen sind
Einermengen]{%
  \begin{aligned}[t]
    &\SemigroupAlg{A}{\star}\dsep
      \MonoidAlg{\PowNE{A}}{\star_{\mathcal P}}{E}\vdash{}\\[-2pt]
    &\exists e\in A\,\bigl(E=\{e\}\land\MonoidAlg{A}{\star}{e}\bigr)
  \end{aligned}%
}{PowerSemigroupIdentityCharacterization}{%
  \DeltaRow{Mengen}{A\dsep E\dsep e\dsep u\dsep v\dsep x}
  \DeltaRow{Funktionen}{\star\dsep\star_{\mathcal P}}
}
\begin{proof}
\par\medskip\noindent\textbf{(i) Die neutrale Menge ist eine Einermenge.}\par
Da \(E\in\PowNE A\), ist \(E\) nichtleer.  Seien \(u,v\in E\).  Aus der
Neutralität von \(E\) folgen
\[
  \{u\}\star_{\mathcal P}E=\{u\},
  \qquad
  E\star_{\mathcal P}\{u\}=\{u\},
  \qquad
  \{v\}\star_{\mathcal P}E=\{v\};
\]
also \(u\star v=u\), \(v\star u=u\) und \(v\star u=v\).  Somit ist
\(u=v\).  Folglich ist
\(E=\{e\}\) für ein \(e\in A\).  \par\medskip\noindent\textbf{(ii) Die Grundhalbgruppe ist ein Monoid.}\par
Für jedes \(x\in A\) liefern nun
\[
  \{x\}\star_{\mathcal P}\{e\}=\{x\}
  \quad\text{und}\quad
  \{e\}\star_{\mathcal P}\{x\}=\{x\}
\]
die Gleichungen \(x\star e=x=e\star x\).  Zusammen mit der vorausgesetzten
Halbgruppenstruktur ist \((A,\star,e)\) daher ein Monoid.
\end{proof}

\section{Monoide als Produktquadrate}

\ProofWideReasonsNearFormula

Wie in Band~28 bezeichnet \(A^2\) das Produktquadrat
\(A\star_{\mathcal P}A\). Auf Teilträgern verwenden wir die
eingeschränkten Operationen. Alle folgenden Aussagen über Quadrate
setzen \(A\neq\varnothing\) voraus.

\FormulaThmDeltaK[Einermengen im nichtleeren Potenzträger]{%
  a\in A\vdash\{a\}\in\PowNE A
}{MonoidPowerSingletonCarrier}{%
  \DeltaRow{Mengen}{A\dsep a}
}
\begin{tabproofwide}
  \proofstepwidestarbreakkeep[1]{a\in A}{\rA}
  \proofstepwidestarbreakkeep[1]{\{a\}\subseteq A}{%
    \FormulaRefAuto{a\in A\vdash\{a\}\subseteq A}{1}}
  \proofstepwidestarbreakkeep[]{a\in\{a\}}{\FormulaRefAuto{a\in\{a\}}}
  \proofstepwidestarbreakkeep[]{\{a\}\neq\varnothing}{%
    \FormulaRefAuto{A\neq\varnothing\eqvdash\exists x\,(x\in A)}{\rEI{3}}}
  \proofstepwidestarbreakkeep[1]{\{a\}\in\PowNE A}{%
    \FormulaRefAuto{NonemptyPowersetMembership}[thm]{\rAI{2,4}}}
\end{tabproofwide}

\FormulaThmDeltaK[Ein Element eines einelementigen Mengenprodukts]{%
  \begin{aligned}[t]
    &\SemigroupAlg{A}{\star}\dsep X,Y\in\PowNE A\dsep
      X\star_{\mathcal P}Y=\{z\}\dsep x\in X\dsep y\in Y\vdash{}\\[-2pt]
    &x\star y=z
  \end{aligned}
}{MonoidPowerSingletonProductValue}{%
  \DeltaRow{Mengen}{A\dsep X\dsep Y\dsep x\dsep y\dsep z}
  \DeltaRow{Funktionen}{\star\dsep\star_{\mathcal P}}
}
\begin{tabproofwide}
  \proofstepwidestarbreakkeep[1]{\SemigroupAlg{A}{\star}}{\rA}
  \proofstepwidestarbreakkeep[2]{X,Y\in\PowNE A}{\rA}
  \proofstepwidestarbreakkeep[3]{X\star_{\mathcal P}Y=\{z\}}{\rA}
  \proofstepwidestarbreakkeep[4]{x\in X\dsep y\in Y}{\rA}
  \proofstepwidestarbreakkeep[1,2,4]{x\star y\in X\star_{\mathcal P}Y}{%
    \FormulaRefAuto{PowerSemigroupProductContains}[thm]{1,2,4}}
  \proofstepwidestarbreakkeep[1,2,3,4]{x\star y\in\{z\}}{\rIE{3,5}}
  \proofstepwidestarbreakkeep[1,2,3,4]{x\star y=z}{%
    \FormulaRefAuto{x\in\{a\}\eqvdash x=a}{6}}
\end{tabproofwide}

\FormulaThmDeltaK[Monoidales Quadrat und sein Potenzträger]{%
  \begin{aligned}[t]
    &\SemigroupAlg{A}{\star}\dsep A\neq\varnothing\dsep
      \MonoidAlg{A^2}{\star}{e}\vdash{}\\[-2pt]
    &(\PowNE A)^2=\PowNE{A^2}\ \land\
      \MonoidAlg{(\PowNE A)^2}{\star_{\mathcal P}}{\{e\}}
  \end{aligned}
}{PowerSemigroupSquareMonoidForward}{%
  \DeltaRow{Mengen}{A\dsep e\dsep X}
  \DeltaRow{Funktionen}{\star\dsep\star_{\mathcal P}}
}
\begin{tabproofsplitwide}
  \proofpartwide{Potenzmonoid und erste Trägerinklusion}
  \proofstepwidestarbreakkeep[1]{\SemigroupAlg{A}{\star}}{\rA}
  \proofstepwidestarbreakkeep[2]{A\neq\varnothing}{\rA}
  \proofstepwidestarbreakkeep[3]{\MonoidAlg{A^2}{\star}{e}}{\rA}
  \proofstepwidestarbreakkeep[1,2]{A^2\subseteq A}{%
    \FormulaRefAuto{SemigroupSquareSubset}[thm]{1,2}}
  \proofstepwidestarbreakkeep[3]{\MonoidAlg{\PowNE{A^2}}{\star_{\mathcal P}}{\{e\}}}{%
    \FormulaRefAuto{PowerSemigroupMonoid}[thm]{3}}
  \proofstepwidestarbreakkeep[1,2]{(\PowNE A)^2\subseteq\PowNE{A^2}}{%
    \FormulaRefAuto{PowerSemigroupSquareSubset}[thm]{1,2}}
  \proofstepwidestarbreakkeep[3]{e\in A^2}{%
    \FormulaRefAuto{MonoidIdentityCarrierAxiom}[ax]{3}}
  \proofstepwidestarbreakkeep[1,2,3]{e\in A}{%
    \FormulaRefAuto{A\subseteq B,\,x\in A\vdash x\in B}{4,7}}
  \proofstepwidestarbreakkeep[1,2,3]{\{e\}\in\PowNE A}{%
    \FormulaRefAuto{MonoidPowerSingletonCarrier}[thm]{8}}
  \proofstepwidestarbreakkeep[1,2,3]{\PowNE A\neq\varnothing}{%
    \FormulaRefAuto{A\neq\varnothing\eqvdash\exists x\,(x\in A)}{\rEI{9}}}
  \proofstepwidestarbreakkeep[1]{\SemigroupAlg{\PowNE A}{\star_{\mathcal P}}}{%
    \FormulaRefAuto{NonemptyPowerSemigroup}[thm]{1}}
  \closeproofpartwide

  \proofpartwide{Umgekehrte Trägerinklusion}
  \setcounter{proofstepnr}{11}
  \proofstepwidestarbreakkeep[12]{X\in\PowNE{A^2}}{\rA}
  \proofstepwidestar[1,2,12]{\PowNE{A^2}\subseteq\PowNE A}{%
      \FormulaRefAuto{NonemptyPowersetMonotone}[thm]{4}}
  \proofstepwidestarbreakkeep[1,2,12]{X\in\PowNE A}{%
    \FormulaRefAuto{A\subseteq B,\,x\in A\vdash x\in B}{%
      13,12}}
  \proofstepwidestarbreakkeep[3,12]{\{e\}\star_{\mathcal P}X=X}{%
    \FormulaRefAuto{MonoidLeftIdentityAxiom}[ax]{5,12}}
  \proofstepwidestarbreakkeep[1,2,3,12]{\{e\}\star_{\mathcal P}X\in(\PowNE A)^2}{%
    \FormulaRefAuto{SemigroupSquareProductMember}[thm]{11,10,9,14}}
  \proofstepwidestarbreakkeep[1,2,3,12]{X\in(\PowNE A)^2}{\rIE{15,16}}
  \proofstepwidestarbreakkeep[1,2,3]{\PowNE{A^2}\subseteq(\PowNE A)^2}{%
    \FormulaRefAuto{SubsetIntroductionRule}[thm]{[12]\vdots 17}}
  \proofstepwidestarbreakkeep[1,2,3]{(\PowNE A)^2=\PowNE{A^2}}{%
    \FormulaRefAuto{A\subseteq B,\,B\subseteq A\vdash A=B}{6,18}}
  \closeproofpartwide

  \proofpartwide{Transport der Monoidstruktur}
  \setcounter{proofstepnr}{19}
  \proofstepwidestarbreakkeep[1,2,3]{\MonoidAlg{(\PowNE A)^2}{\star_{\mathcal P}}{\{e\}}}{\rIE{19,5}}
  \proofstepwidestarbreakkeep[1,2,3]{%
    (\PowNE A)^2=\PowNE{A^2}\ \land\
    \MonoidAlg{(\PowNE A)^2}{\star_{\mathcal P}}{\{e\}}}{\rAI{19,20}}
  \closeproofpartwide
\end{tabproofsplitwide}

\FormulaThmDeltaK[Elemente der neutralen Menge wirken neutral auf dem Quadrat]{%
  \begin{aligned}[t]
    &\SemigroupAlg{A}{\star}\dsep A\neq\varnothing\dsep
      \MonoidAlg{(\PowNE A)^2}{\star_{\mathcal P}}{E}\dsep
      e\in E\dsep t\in A^2\vdash{}\\[-2pt]
    &e\in A^2\ \land\ (e\star t=t\land t\star e=t)
  \end{aligned}
}{PowerSemigroupSquareIdentityActions}{%
  \DeltaRow{Mengen}{A\dsep E\dsep e\dsep t}
  \DeltaRow{Funktionen}{\star\dsep\star_{\mathcal P}}
}
\begin{tabproofsplitwide}
  \proofpartwide{Trägerzugehörigkeit der neutralen Elemente}
  \proofstepwidestarbreakkeep[1]{\SemigroupAlg{A}{\star}}{\rA}
  \proofstepwidestarbreakkeep[2]{A\neq\varnothing}{\rA}
  \proofstepwidestarbreakkeep[3]{\MonoidAlg{(\PowNE A)^2}{\star_{\mathcal P}}{E}}{\rA}
  \proofstepwidestarbreakkeep[4]{e\in E}{\rA}
  \proofstepwidestarbreakkeep[5]{t\in A^2}{\rA}
  \proofstepwidestarbreakkeep[3]{E\in(\PowNE A)^2}{%
    \FormulaRefAuto{MonoidIdentityCarrierAxiom}[ax]{3}}
  \proofstepwidestarbreakkeep[1,2]{(\PowNE A)^2\subseteq\PowNE{A^2}}{%
    \FormulaRefAuto{PowerSemigroupSquareSubset}[thm]{1,2}}
  \proofstepwidestarbreakkeep[1,2,3]{E\in\PowNE{A^2}}{%
    \FormulaRefAuto{A\subseteq B,\,x\in A\vdash x\in B}{7,6}}
  \proofstepwidestarbreakkeep[1,2,3,4]{e\in A^2}{%
    \FormulaRefAuto{PowerSemigroupCarrierElement}[thm]{8,4}}
  \closeproofpartwide

  \proofpartwide{Links- und Rechtsneutralität auf dem Quadrat}
  \setcounter{proofstepnr}{9}
  \proofstepwidestarbreakkeep[1,2]{A^2\subseteq A}{%
    \FormulaRefAuto{SemigroupSquareSubset}[thm]{1,2}}
  \proofstepwidestarbreakkeep[1,2,5]{t\in A}{%
    \FormulaRefAuto{A\subseteq B,\,x\in A\vdash x\in B}{10,5}}
  \proofstepwidestarbreakkeep[1,2,5]{\{t\}\in(\PowNE A)^2}{%
    \FormulaRefAuto{PowerSemigroupSquareSingletonCriterion}[thm]{1,2,11,5}}
  \proofstepwidestarbreakkeep[1,2,3,5]{E\star_{\mathcal P}\{t\}=\{t\}}{%
    \FormulaRefAuto{MonoidLeftIdentityAxiom}[ax]{3,12}}
  \proofstepwidestarbreakkeep[1,2,3,5]{\{t\}\star_{\mathcal P}E=\{t\}}{%
    \FormulaRefAuto{MonoidRightIdentityAxiom}[ax]{3,12}}
  \proofstepwidestar[1,2,3]{\PowNE{A^2}\subseteq\PowNE A}{%
      \FormulaRefAuto{NonemptyPowersetMonotone}[thm]{10}}
  \proofstepwidestarbreakkeep[1,2,3]{E\in\PowNE A}{%
    \FormulaRefAuto{A\subseteq B,\,x\in A\vdash x\in B}{%
      15,8}}
  \proofstepwidestarbreakkeep[1,2,5]{\{t\}\in\PowNE A}{%
    \FormulaRefAuto{MonoidPowerSingletonCarrier}[thm]{11}}
  \proofstepwidestarbreakkeep[]{t\in\{t\}}{\FormulaRefAuto{a\in\{a\}}}
  \proofstepwidestarbreakkeep[1,2,3,4,5]{e\star t=t}{%
    \FormulaRefAuto{MonoidPowerSingletonProductValue}[thm]{1,16,17,13,4,18}}
  \proofstepwidestarbreakkeep[1,2,3,4,5]{t\star e=t}{%
    \FormulaRefAuto{MonoidPowerSingletonProductValue}[thm]{1,17,16,14,18,4}}
  \proofstepwidestarbreakkeep[1,2,3,4,5]{e\in A^2\ \land\ (e\star t=t\land t\star e=t)}{%
    \rAI{9,\rAI{19,20}}}
  \closeproofpartwide
\end{tabproofsplitwide}

\FormulaThmDeltaK[Monoidstruktur des Potenzquadrats erkennt das Grundquadrat]{%
  \begin{aligned}[t]
    &\SemigroupAlg{A}{\star}\dsep A\neq\varnothing\dsep
      \MonoidAlg{(\PowNE A)^2}{\star_{\mathcal P}}{E}\vdash{}\\[-2pt]
    &\exists e\in A^2\,\bigl(E=\{e\}\land\MonoidAlg{A^2}{\star}{e}\bigr)
  \end{aligned}
}{PowerSemigroupSquareMonoidBackward}{%
  \DeltaRow{Mengen}{A\dsep E\dsep e\dsep t\dsep v}
  \DeltaRow{Funktionen}{\star\dsep\star_{\mathcal P}}
}
\begin{tabproofsplitwide}
  \proofpartwide{Einermengigkeit des neutralen Mengenprodukts}
  \proofstepwidestarbreakkeep[1]{\SemigroupAlg{A}{\star}}{\rA}
  \proofstepwidestarbreakkeep[2]{A\neq\varnothing}{\rA}
  \proofstepwidestarbreakkeep[3]{\MonoidAlg{(\PowNE A)^2}{\star_{\mathcal P}}{E}}{\rA}
  \proofstepwidestarbreakkeep[3]{E\in(\PowNE A)^2}{%
    \FormulaRefAuto{MonoidIdentityCarrierAxiom}[ax]{3}}
  \proofstepwidestarbreakkeep[1,2]{(\PowNE A)^2\subseteq\PowNE{A^2}}{%
    \FormulaRefAuto{PowerSemigroupSquareSubset}[thm]{1,2}}
  \proofstepwidestarbreakkeep[1,2,3]{E\in\PowNE{A^2}}{%
    \FormulaRefAuto{A\subseteq B,\,x\in A\vdash x\in B}{5,4}}
  \proofstepwidestar[1,2,3]{E\subseteq A^2\land E\neq\varnothing}{%
      \FormulaRefAuto{NonemptyPowersetMembership}[thm]{6}}
  \proofstepwidestarbreakkeep[1,2,3]{E\neq\varnothing}{%
    \rAEb{7}}
  \proofstepwidestarbreakkeep[1,2,3]{\exists e\,(e\in E)}{%
    \FormulaRefAuto{A\neq\varnothing\eqvdash\exists x\,(x\in A)}{8}}
  \proofstepwidestarbreakkeep[10]{e\in E}{\rA}
  \proofstepwidestarbreakkeep[1,2,3,10]{e\in A^2}{%
    \FormulaRefAuto{PowerSemigroupCarrierElement}[thm]{6,10}}
  \proofstepwidestarbreakkeep[12]{v\in E}{\rA}
  \proofstepwidestarbreakkeep[1,2,3,12]{v\in A^2}{%
    \FormulaRefAuto{PowerSemigroupCarrierElement}[thm]{6,12}}
  \proofstepwidestar[1,2,3,10,12]{e\in A^2\land(e\star v=v\land v\star e=v)}{%
      \FormulaRefAuto{PowerSemigroupSquareIdentityActions}[thm]{1,2,3,10,13}}
  \proofstepwidestarbreakkeep[1,2,3,10,12]{e\star v=v}{%
    \rAEa{\rAEb{14}}}
  \proofstepwidestar[1,2,3,10,12]{v\in A^2\land(v\star e=e\land e\star v=e)}{%
      \FormulaRefAuto{PowerSemigroupSquareIdentityActions}[thm]{1,2,3,12,11}}
  \proofstepwidestarbreakkeep[1,2,3,10,12]{e\star v=e}{%
    \rAEb{\rAEb{16}}}
  \proofstepwidestarbreakkeep[1,2,3,10,12]{v=e}{\rIE{15,17}}
  \proofstepwidestarbreakkeep[19]{v=e}{\rA}
  \proofstepwidestarbreakkeep[10,19]{v\in E}{\rIE{19,10}}
  \proofstepwidestarbreakkeep[1,2,3,10]{\forall v\,(v\in E\leftrightarrow v=e)}{%
    \rUI{\rLRI{\rRI{12,18},\rRI{19,20}}}}
  \proofstepwidestar[1,2,3,10]{v\in\{e\}\leftrightarrow v=e}{%
      \FormulaRefAuto{x\in\{a\}\eqvdash x=a}}
  \proofstepwidestarbreakkeep[1,2,3,10]{E=\{e\}}{%
    \FormulaRefAuto{\forall x\,(x\in A\leftrightarrow x\in B)\eqvdash A=B}{%
      21,22}}
  \closeproofpartwide

  \proofpartwide{Monoidstruktur des Grundquadrats}
  \setcounter{proofstepnr}{23}
  \proofstepwidestarbreakkeep[1,2]{\mathsf{UHGr}(A^2;A,\star)}{%
    \FormulaRefAuto{SemigroupSquareIsSubsemigroup}[thm]{1,2}}
  \proofstepwidestarbreakkeep[1,2]{\SemigroupAlg{A^2}{\star}}{%
    \FormulaRefAuto{SubsemigroupFormsSemigroup}[thm]{1,24}}
  \proofstepwidestarbreakkeep[26]{t\in A^2}{\rA}
  \proofstepwidestar[1,2,3,10,26]{e\in A^2\land(e\star t=t\land t\star e=t)}{%
      \FormulaRefAuto{PowerSemigroupSquareIdentityActions}[thm]{1,2,3,10,26}}
  \proofstepwidestarbreakkeep[1,2,3,10,26]{e\star t=t\land t\star e=t}{%
    \rAEb{27}}
  \proofstepwidestarbreakkeep[1,2,3,10]{\MonoidAlg{A^2}{\star}{e}}{%
    \FormulaRefAuto{MonoidDef}[def]{25,11,\rAEa{28},\rAEb{28}}}
  \proofstepwidestarbreakkeep[1,2,3,10]{\exists e\in A^2\,(E=\{e\}\land\MonoidAlg{A^2}{\star}{e})}{%
    \rEI{11,\rAI{23,29}}}
  \proofstepwidestarbreakkeep[1,2,3]{\exists e\in A^2\,(E=\{e\}\land\MonoidAlg{A^2}{\star}{e})}{%
    \rEE{9,10,30}}
  \closeproofpartwide
\end{tabproofsplitwide}

\FormulaThmDeltaK[Monoidkriterium für Produktquadrate]{%
  \begin{aligned}[t]
    &\SemigroupAlg{A}{\star}\dsep A\neq\varnothing\vdash{}\\[-2pt]
    &\bigl(\exists E\,\MonoidAlg{(\PowNE A)^2}{\star_{\mathcal P}}{E}\bigr)
      \leftrightarrow\bigl(\exists e\,\MonoidAlg{A^2}{\star}{e}\bigr)
  \end{aligned}
}{PowerSemigroupSquareMonoidCriterion}{%
  \DeltaRow{Mengen}{A\dsep E\dsep e}
  \DeltaRow{Funktionen}{\star\dsep\star_{\mathcal P}}
}
\begin{tabproofsplitwide}
  \proofpartwide{Vom Potenzquadrat zum Grundquadrat}
  \proofstepwidestarbreakkeep[1]{\SemigroupAlg{A}{\star}}{\rA}
  \proofstepwidestarbreakkeep[2]{A\neq\varnothing}{\rA}
  \proofstepwidestarbreakkeep[3]{\exists E\,\MonoidAlg{(\PowNE A)^2}{\star_{\mathcal P}}{E}}{\rA}
  \proofstepwidestarbreakkeep[4]{\MonoidAlg{(\PowNE A)^2}{\star_{\mathcal P}}{E}}{\rA}
  \proofstepwidestarbreakkeep[1,2,4]{\exists e\in A^2\,(E=\{e\}\land\MonoidAlg{A^2}{\star}{e})}{%
    \FormulaRefAuto{PowerSemigroupSquareMonoidBackward}[thm]{1,2,4}}
  \proofstepwidestarbreakkeep[6]{e\in A^2\land(E=\{e\}\land\MonoidAlg{A^2}{\star}{e})}{\rA}
  \proofstepwidestarbreakkeep[6]{\MonoidAlg{A^2}{\star}{e}}{\rAEb{\rAEb{6}}}
  \proofstepwidestarbreakkeep[6]{\exists e\,\MonoidAlg{A^2}{\star}{e}}{\rEI{7}}
  \proofstepwidestarbreakkeep[1,2,4]{\exists e\,\MonoidAlg{A^2}{\star}{e}}{\rEE{5,6,8}}
  \proofstepwidestarbreakkeep[1,2,3]{\exists e\,\MonoidAlg{A^2}{\star}{e}}{\rEE{3,4,9}}
  \closeproofpartwide

  \proofpartwide{Vom Grundquadrat zum Potenzquadrat}
  \setcounter{proofstepnr}{10}
  \proofstepwidestarbreakkeep[11]{\exists e\,\MonoidAlg{A^2}{\star}{e}}{\rA}
  \proofstepwidestarbreakkeep[12]{\MonoidAlg{A^2}{\star}{e}}{\rA}
  \proofstepwidestar[1,2,12]{(\PowNE A)^2=\PowNE{A^2}\land\MonoidAlg{(\PowNE A)^2}{\star_{\mathcal P}}{\{e\}}}{%
      \FormulaRefAuto{PowerSemigroupSquareMonoidForward}[thm]{1,2,12}}
  \proofstepwidestarbreakkeep[1,2,12]{\MonoidAlg{(\PowNE A)^2}{\star_{\mathcal P}}{\{e\}}}{%
    \rAEb{13}}
  \proofstepwidestarbreakkeep[1,2,12]{\exists E\,\MonoidAlg{(\PowNE A)^2}{\star_{\mathcal P}}{E}}{\rEI{14}}
  \proofstepwidestarbreakkeep[1,2,11]{\exists E\,\MonoidAlg{(\PowNE A)^2}{\star_{\mathcal P}}{E}}{\rEE{11,12,15}}
  \proofstepwidestarbreakkeep[1,2]{%
    \bigl(\exists E\,\MonoidAlg{(\PowNE A)^2}{\star_{\mathcal P}}{E}\bigr)
    \leftrightarrow\bigl(\exists e\,\MonoidAlg{A^2}{\star}{e}\bigr)}{%
    \rLRI{\rRI{3,10},\rRI{11,16}}}
  \closeproofpartwide
\end{tabproofsplitwide}

\subsection{Der Rückzug auf ein monoidales Quadrat}

\FormulaThmDeltaK[Träger der Produkte mit der Quadrateins]{%
  \begin{aligned}[t]
    &\SemigroupAlg{A}{\star}\dsep A\neq\varnothing\dsep
      \MonoidAlg{A^2}{\star}{e}\dsep a\in A\vdash{}\\[-2pt]
    &e\in A\ \land\ (e\star a\in A^2\land a\star e\in A^2)
  \end{aligned}
}{MonoidSquareIdentityProductCarrier}{%
  \DeltaRow{Mengen}{A\dsep e\dsep a}
  \DeltaRow{Funktionen}{\star}
}
\begin{tabproofwide}
  \proofstepwidestarbreakkeep[1]{\SemigroupAlg{A}{\star}}{\rA}
  \proofstepwidestarbreakkeep[2]{A\neq\varnothing}{\rA}
  \proofstepwidestarbreakkeep[3]{\MonoidAlg{A^2}{\star}{e}}{\rA}
  \proofstepwidestarbreakkeep[4]{a\in A}{\rA}
  \proofstepwidestarbreakkeep[3]{e\in A^2}{%
    \FormulaRefAuto{MonoidIdentityCarrierAxiom}[ax]{3}}
  \proofstepwidestarbreakkeep[1,2]{A^2\subseteq A}{%
    \FormulaRefAuto{SemigroupSquareSubset}[thm]{1,2}}
  \proofstepwidestarbreakkeep[1,2,3]{e\in A}{%
    \FormulaRefAuto{A\subseteq B,\,x\in A\vdash x\in B}{6,5}}
  \proofstepwidestarbreakkeep[1,2,3,4]{e\star a\in A^2}{%
    \FormulaRefAuto{SemigroupSquareProductMember}[thm]{1,2,7,4}}
  \proofstepwidestarbreakkeep[1,2,3,4]{a\star e\in A^2}{%
    \FormulaRefAuto{SemigroupSquareProductMember}[thm]{1,2,4,7}}
  \proofstepwidestarbreakkeep[1,2,3,4]{e\in A\ \land\ (e\star a\in A^2\land a\star e\in A^2)}{%
    \rAI{7,\rAI{8,9}}}
\end{tabproofwide}

\par\noindent\begin{minipage}{\linewidth}
\FormulaThmDeltaK[Die Quadrateins vertauscht mit allen Elementen]{%
  \SemigroupAlg{A}{\star}\dsep A\neq\varnothing\dsep
  \MonoidAlg{A^2}{\star}{e}\dsep a\in A
  \vdash e\star a=a\star e
}{MonoidSquareCentralIdentity}{%
  \DeltaRow{Mengen}{A\dsep e\dsep a}
  \DeltaRow{Funktionen}{\star}
}
\end{minipage}\par
\begin{tabproofwide}
  \proofstepwidestarbreakkeep[1]{\SemigroupAlg{A}{\star}}{\rA}
  \proofstepwidestarbreakkeep[2]{A\neq\varnothing}{\rA}
  \proofstepwidestarbreakkeep[3]{\MonoidAlg{A^2}{\star}{e}}{\rA}
  \proofstepwidestarbreakkeep[4]{a\in A}{\rA}
  \proofstepwidestarbreakkeep[1,2,3,4]{e\in A\ \land\ (e\star a\in A^2\land a\star e\in A^2)}{%
    \FormulaRefAuto{MonoidSquareIdentityProductCarrier}[thm]{1,2,3,4}}
  \proofstepwidestar[1,2,3,4]{(e\star a)\star e=e\star a}{%
      \FormulaRefAuto{MonoidRightIdentityAxiom}[ax]{3,\rAEa{\rAEb{5}}}}
  \proofstepwidestarbreakkeep[1,2,3,4]{e\star a=(e\star a)\star e}{%
    \FormulaRefAuto{a=b\vdash b=a}{%
      6}}
  \proofstepwidestarbreakkeep[1,2,3,4]{=e\star(a\star e)}{%
    \FormulaRefAuto{SemigroupAssociativityAxiom}[ax]{1,\rAEa{5},4}}
  \proofstepwidestarbreakkeep[1,2,3,4]{=a\star e}{%
    \FormulaRefAuto{MonoidLeftIdentityAxiom}[ax]{3,\rAEb{\rAEb{5}}}}
  \proofstepwidestarbreakkeep[1,2,3,4]{e\star a=a\star e}{\rChain{7,8,9}}
\end{tabproofwide}

\FormulaThmDeltaK[Ein Produkt hängt nur von den Rückzugswerten ab]{%
  \SemigroupAlg{A}{\star}\dsep A\neq\varnothing\dsep
  \MonoidAlg{A^2}{\star}{e}\dsep a,b\in A
  \vdash (e\star a)\star(e\star b)=a\star b
}{MonoidSquareProductFactorization}{%
  \DeltaRow{Mengen}{A\dsep e\dsep a\dsep b}
  \DeltaRow{Funktionen}{\star}
}
\begin{tabproofwide}
  \proofstepwidestarbreakkeep[1]{\SemigroupAlg{A}{\star}}{\rA}
  \proofstepwidestarbreakkeep[2]{A\neq\varnothing}{\rA}
  \proofstepwidestarbreakkeep[3]{\MonoidAlg{A^2}{\star}{e}}{\rA}
  \proofstepwidestarbreakkeep[4]{a,b\in A}{\rA}
  \proofstepwidestarbreakkeep[1,2,3,4]{e\in A\ \land\ (e\star a\in A^2\land a\star e\in A^2)}{%
    \FormulaRefAuto{MonoidSquareIdentityProductCarrier}[thm]{1,2,3,4}}
  \proofstepwidestarbreakkeep[1,2,3,4]{e\star a\in A}{%
    \FormulaRefAuto{SemigroupClosureAxiom}[ax]{1,\rAEa{5},4}}
  \proofstepwidestarbreakkeep[1,2,4]{a\star b\in A^2}{%
    \FormulaRefAuto{SemigroupSquareProductMember}[thm]{1,2,4}}
  \proofstepwidestar[1,2,3,4]{((e\star a)\star e)\star b=(e\star a)\star(e\star b)}{%
      \FormulaRefAuto{SemigroupAssociativityAxiom}[ax]{1,6,\rAEa{5},4}}
  \proofstepwidestarbreakkeep[1,2,3,4]{(e\star a)\star(e\star b)=((e\star a)\star e)\star b}{%
    \FormulaRefAuto{a=b\vdash b=a}{%
      8}}
  \proofstepwidestar[1,2,3,4]{(e\star a)\star e=e\star a}{%
      \FormulaRefAuto{MonoidRightIdentityAxiom}[ax]{3,\rAEa{\rAEb{5}}}}
  \proofstepwidestarbreakkeep[1,2,3,4]{=(e\star a)\star b}{%
    \rIE{10,9}}
  \proofstepwidestarbreakkeep[1,2,3,4]{=e\star(a\star b)}{%
    \FormulaRefAuto{SemigroupAssociativityAxiom}[ax]{1,\rAEa{5},4}}
  \proofstepwidestarbreakkeep[1,2,3,4]{=a\star b}{%
    \FormulaRefAuto{MonoidLeftIdentityAxiom}[ax]{3,7}}
  \proofstepwidestarbreakkeep[1,2,3,4]{(e\star a)\star(e\star b)=a\star b}{\rChain{9,11,12,13}}
\end{tabproofwide}

\FormulaDefDeltaK[Rückzug auf das monoidale Quadrat]{%
  \begin{aligned}[t]
    &\SemigroupAlg{A}{\star}\dsep A\neq\varnothing\dsep
      \MonoidAlg{A^2}{\star}{e}\vdash{}\\[-2pt]
    &r_e\colon A\to A^2,\qquad
      \forall a\in A\,\bigl(r_e(a)\coloneqq e\star a\bigr)
  \end{aligned}
}{MonoidSquareInflationDef}{%
  \DeltaRow{Mengen}{A\dsep e\dsep a}
  \DeltaRow{Funktionen}{\star\dsep r_e}
}
\begin{tabproofwide}
  \proofstepwidestarbreakkeep[1]{\SemigroupAlg{A}{\star}}{\rA}
  \proofstepwidestarbreakkeep[2]{A\neq\varnothing}{\rA}
  \proofstepwidestarbreakkeep[3]{\MonoidAlg{A^2}{\star}{e}}{\rA}
  \proofstepwidestarbreakkeep[4]{a\in A}{\rA}
  \proofstepwidestar[1,2,3,4]{e\in A\land(e\star a\in A^2\land a\star e\in A^2)}{%
      \FormulaRefAuto{MonoidSquareIdentityProductCarrier}[thm]{1,2,3,4}}
  \proofstepwidestarbreakkeep[1,2,3,4]{e\star a\in A^2}{%
    \rAEa{\rAEb{5}}}
  \proofstepwidestarbreakkeep[1,2,3]{\forall a\in A\,(e\star a\in A^2)}{\rUI{\rRI{4,6}}}
  \proofstepwidestarbreakkeep[1,2,3]{\exists!r\,(r\colon A\to A^2\land\forall a\in A\,(r(a)=e\star a))}{%
    \FormulaRefAuto{TypedFunctionDefinitionPrinciple}[thm]{7}}
\end{tabproofwide}

\FormulaThmDeltaK[Der Rückzug fixiert das Quadrat]{%
  \SemigroupAlg{A}{\star}\dsep A\neq\varnothing\dsep
  \MonoidAlg{A^2}{\star}{e}\dsep g\in A^2
  \vdash r_e(g)=g
}{MonoidSquareRetractionValue}{%
  \DeltaRow{Mengen}{A\dsep e\dsep g}
  \DeltaRow{Funktionen}{\star\dsep r_e}
}
\begin{tabproofwide}
  \proofstepwidestarbreakkeep[1]{\SemigroupAlg{A}{\star}}{\rA}
  \proofstepwidestarbreakkeep[2]{A\neq\varnothing}{\rA}
  \proofstepwidestarbreakkeep[3]{\MonoidAlg{A^2}{\star}{e}}{\rA}
  \proofstepwidestarbreakkeep[4]{g\in A^2}{\rA}
  \proofstepwidestar[1,2,4]{A^2\subseteq A}{%
      \FormulaRefAuto{SemigroupSquareSubset}[thm]{1,2}}
  \proofstepwidestarbreakkeep[1,2,4]{g\in A}{%
    \FormulaRefAuto{A\subseteq B,\,x\in A\vdash x\in B}{%
      5,4}}
  \proofstepwidestarbreakkeep[1,2,3,4]{r_e(g)=e\star g}{%
    \FormulaRefAuto{MonoidSquareInflationDef}[def]{1,2,3,6}}
  \proofstepwidestarbreakkeep[3,4]{=g}{%
    \FormulaRefAuto{MonoidLeftIdentityAxiom}[ax]{3,4}}
  \proofstepwidestarbreakkeep[1,2,3,4]{r_e(g)=g}{\rChain{7,8}}
\end{tabproofwide}

\FormulaThmDeltaK[Produktformel des Rückzugs]{%
  \SemigroupAlg{A}{\star}\dsep A\neq\varnothing\dsep
  \MonoidAlg{A^2}{\star}{e}\dsep a,b\in A
  \vdash a\star b=r_e(a)\star r_e(b)
}{MonoidSquareInflationProduct}{%
  \DeltaRow{Mengen}{A\dsep e\dsep a\dsep b}
  \DeltaRow{Funktionen}{\star\dsep r_e}
}
\begin{tabproofwide}
  \proofstepwidestarbreakkeep[1]{\SemigroupAlg{A}{\star}}{\rA}
  \proofstepwidestarbreakkeep[2]{A\neq\varnothing}{\rA}
  \proofstepwidestarbreakkeep[3]{\MonoidAlg{A^2}{\star}{e}}{\rA}
  \proofstepwidestarbreakkeep[4]{a,b\in A}{\rA}
  \proofstepwidestarbreakkeep[1,2,3,4]{r_e(a)=e\star a\dsep r_e(b)=e\star b}{%
    \FormulaRefAuto{MonoidSquareInflationDef}[def]{1,2,3,4}}
  \proofstepwidestarbreakkeep[1,2,3,4]{(e\star a)\star(e\star b)=a\star b}{%
    \FormulaRefAuto{MonoidSquareProductFactorization}[thm]{1,2,3,4}}
  \proofstepwidestarbreakkeep[1,2,3,4]{a\star b=r_e(a)\star r_e(b)}{%
    \rIE{5,\FormulaRefAuto{a=b\vdash b=a}{6}}}
\end{tabproofwide}

\FormulaThmDeltaK[Eine Halbgruppe mit monoidalem Quadrat ist dessen Inflation]{%
  \begin{aligned}[t]
    &\SemigroupAlg{A}{\star}\dsep A\neq\varnothing\dsep
      \MonoidAlg{A^2}{\star}{e}\vdash{}\\[-2pt]
    &\mathsf{Infl}_{r_e}\bigl((A^2,\star),(A,\star)\bigr)
  \end{aligned}
}{MonoidSquareInflationProperties}{%
  \DeltaRow{Mengen}{A\dsep e\dsep g\dsep a\dsep b}
  \DeltaRow{Funktionen}{\star\dsep r_e}
}
\begin{tabproofwide}
  \proofstepwidestarbreakkeep[1]{\SemigroupAlg{A}{\star}}{\rA}
  \proofstepwidestarbreakkeep[2]{A\neq\varnothing}{\rA}
  \proofstepwidestarbreakkeep[3]{\MonoidAlg{A^2}{\star}{e}}{\rA}
  \proofstepwidestarbreakkeep[3]{\SemigroupAlg{A^2}{\star}}{%
    \FormulaRefAuto{MonoidBaseSemigroupAxiom}[ax]{3}}
  \proofstepwidestarbreakkeep[1,2]{A^2\subseteq A}{%
    \FormulaRefAuto{SemigroupSquareSubset}[thm]{1,2}}
  \proofstepwidestarbreakkeep[1,2,3]{r_e\colon A\to A^2}{%
    \FormulaRefAuto{MonoidSquareInflationDef}[def]{1,2,3}}
  \proofstepwidestarbreakkeep[7]{g\in A^2}{\rA}
  \proofstepwidestarbreakkeep[1,2,3,7]{r_e(g)=g}{%
    \FormulaRefAuto{MonoidSquareRetractionValue}[thm]{1,2,3,7}}
  \proofstepwidestarbreakkeep[1,2,3]{\forall g\in A^2\,(r_e(g)=g)}{\rUI{\rRI{7,8}}}
  \proofstepwidestarbreakkeep[10]{a,b\in A}{\rA}
  \proofstepwidestarbreakkeep[1,2,3,10]{a\star b=r_e(a)\star r_e(b)}{%
    \FormulaRefAuto{MonoidSquareInflationProduct}[thm]{1,2,3,10}}
  \proofstepwidestarbreakkeep[1,2,3]{\forall a,b\in A\,(a\star b=r_e(a)\star r_e(b))}{\rUI{\rUI{\rRI{10,11}}}}
  \proofstepwidestarbreakkeep[1,2,3]{\mathsf{Infl}_{r_e}\bigl((A^2,\star),(A,\star)\bigr)}{%
    \FormulaRefAuto{SemigroupInflationDef}[def]{4,1,5,6,9,12}}
\end{tabproofwide}

\par\noindent\begin{minipage}{\linewidth}
\FormulaThmDeltaK[Ein einelementiges Rückzugsbild bestimmt jeden Rückzugswert]{%
  \begin{aligned}[t]
    &\SemigroupAlg{A}{\star}\dsep A\neq\varnothing\dsep
      \MonoidAlg{A^2}{\star}{e}\dsep X\in\PowNE A\dsep{}\\[-2pt]
    &\{e\}\star_{\mathcal P}X=\{g\}\dsep a\in X\vdash{}\\[-2pt]
    &r_e(a)=g
  \end{aligned}
}{MonoidSquarePowerFiberElementForward}{%
  \DeltaRow{Mengen}{A\dsep e\dsep g\dsep X\dsep a}
  \DeltaRow{Funktionen}{\star\dsep\star_{\mathcal P}\dsep r_e}
}
\end{minipage}\par
\begin{tabproofwide}
  \proofstepwidestarbreakkeep[1]{\SemigroupAlg{A}{\star}}{\rA}
  \proofstepwidestarbreakkeep[2]{A\neq\varnothing}{\rA}
  \proofstepwidestarbreakkeep[3]{\MonoidAlg{A^2}{\star}{e}}{\rA}
  \proofstepwidestarbreakkeep[4]{X\in\PowNE A}{\rA}
  \proofstepwidestarbreakkeep[5]{\{e\}\star_{\mathcal P}X=\{g\}}{\rA}
  \proofstepwidestarbreakkeep[6]{a\in X}{\rA}
  \proofstepwidestarbreakkeep[4,6]{a\in A}{%
    \FormulaRefAuto{PowerSemigroupCarrierElement}[thm]{4,6}}
  \proofstepwidestar[1,2,3,4,6]{e\in A\land(e\star a\in A^2\land a\star e\in A^2)}{%
      \FormulaRefAuto{MonoidSquareIdentityProductCarrier}[thm]{1,2,3,7}}
  \proofstepwidestarbreakkeep[1,2,3,4,6]{e\in A}{%
    \rAEa{8}}
  \proofstepwidestarbreakkeep[1,2,3,4,6]{\{e\}\in\PowNE A}{%
    \FormulaRefAuto{MonoidPowerSingletonCarrier}[thm]{9}}
  \proofstepwidestarbreakkeep[]{e\in\{e\}}{\FormulaRefAuto{a\in\{a\}}}
  \proofstepwidestarbreakkeep[1,2,3,4,5,6]{e\star a=g}{%
    \FormulaRefAuto{MonoidPowerSingletonProductValue}[thm]{1,10,4,5,11,6}}
  \proofstepwidestarbreakkeep[1,2,3,4,6]{r_e(a)=e\star a}{%
    \FormulaRefAuto{MonoidSquareInflationDef}[def]{1,2,3,7}}
  \proofstepwidestarbreakkeep[1,2,3,4,5,6]{r_e(a)=g}{\rChain{13,12}}
\end{tabproofwide}

\FormulaThmDeltaK[Konstante Rückzugswerte bestimmen jedes Element des Mengenprodukts]{%
  \begin{aligned}[t]
    &\SemigroupAlg{A}{\star}\dsep A\neq\varnothing\dsep
      \MonoidAlg{A^2}{\star}{e}\dsep X\in\PowNE A\dsep{}\\[-2pt]
    &\forall a\in X\,(r_e(a)=g)\dsep
      z\in\{e\}\star_{\mathcal P}X\vdash z=g
  \end{aligned}
}{MonoidSquarePowerFiberProductElement}{%
  \DeltaRow{Mengen}{A\dsep e\dsep g\dsep X\dsep a\dsep v\dsep z}
  \DeltaRow{Funktionen}{\star\dsep\star_{\mathcal P}\dsep r_e}
}
\begin{tabproofwide}
  \proofstepwidestarbreakkeep[1]{\SemigroupAlg{A}{\star}}{\rA}
  \proofstepwidestarbreakkeep[2]{A\neq\varnothing}{\rA}
  \proofstepwidestarbreakkeep[3]{\MonoidAlg{A^2}{\star}{e}}{\rA}
  \proofstepwidestarbreakkeep[4]{X\in\PowNE A}{\rA}
  \proofstepwidestarbreakkeep[5]{\forall a\in X\,(r_e(a)=g)}{\rA}
  \proofstepwidestarbreakkeep[6]{z\in\{e\}\star_{\mathcal P}X}{\rA}
  \proofstepwidestarbreakkeep[3]{e\in A^2}{\FormulaRefAuto{MonoidIdentityCarrierAxiom}[ax]{3}}
  \proofstepwidestar[1,2,3]{A^2\subseteq A}{%
      \FormulaRefAuto{SemigroupSquareSubset}[thm]{1,2}}
  \proofstepwidestarbreakkeep[1,2,3]{e\in A}{%
    \FormulaRefAuto{A\subseteq B,\,x\in A\vdash x\in B}{%
      8,7}}
  \proofstepwidestarbreakkeep[1,2,3]{\{e\}\in\PowNE A}{%
    \FormulaRefAuto{MonoidPowerSingletonCarrier}[thm]{9}}
  \proofstepwidestarbreakkeep[1,2,3,4,6]{\exists v\in\{e\}\,\exists a\in X\,(z=v\star a)}{%
    \FormulaRefAuto{PowerSemigroupProductWitnesses}[thm]{1,10,4,6}}
  \proofstepwidestarbreakkeep[12]{v\in\{e\}\land(a\in X\land z=v\star a)}{\rA}
  \proofstepwidestarbreakkeep[12]{v=e}{\FormulaRefAuto{x\in\{a\}\eqvdash x=a}{\rAEa{12}}}
  \proofstepwidestarbreakkeep[12]{a\in X}{\rAEa{\rAEb{12}}}
  \proofstepwidestarbreakkeep[4,12]{a\in A}{\FormulaRefAuto{PowerSemigroupCarrierElement}[thm]{4,14}}
  \proofstepwidestarbreakkeep[1,2,3,4,12]{r_e(a)=e\star a}{%
    \FormulaRefAuto{MonoidSquareInflationDef}[def]{1,2,3,15}}
  \proofstepwidestarbreakkeep[5,12]{r_e(a)=g}{\rRE{14,\rUE{5}}}
  \proofstepwidestarbreakkeep[1,2,3,4,5,12]{z=g}{%
    \rIE{13,16,17,\rAEb{\rAEb{12}}}}
  \proofstepwidestarbreakkeep[1,2,3,4,5,6]{z=g}{\rEEStar{11,12\text{--}18}}
\end{tabproofwide}

\FormulaThmDeltaK[Konstante Rückzugswerte liefern ein einelementiges Mengenprodukt]{%
  \begin{aligned}[t]
    &\SemigroupAlg{A}{\star}\dsep A\neq\varnothing\dsep
      \MonoidAlg{A^2}{\star}{e}\dsep X\in\PowNE A\dsep
      \forall a\in X\,(r_e(a)=g)\vdash{}\\[-2pt]
    &\{e\}\star_{\mathcal P}X=\{g\}
  \end{aligned}
}{MonoidSquarePowerFiberBackward}{%
  \DeltaRow{Mengen}{A\dsep e\dsep g\dsep X\dsep a\dsep z}
  \DeltaRow{Funktionen}{\star\dsep\star_{\mathcal P}\dsep r_e}
}
\begin{tabproofwide}
  \proofstepwidestarbreakkeep[1]{\SemigroupAlg{A}{\star}}{\rA}
  \proofstepwidestarbreakkeep[2]{A\neq\varnothing}{\rA}
  \proofstepwidestarbreakkeep[3]{\MonoidAlg{A^2}{\star}{e}}{\rA}
  \proofstepwidestarbreakkeep[4]{X\in\PowNE A}{\rA}
  \proofstepwidestarbreakkeep[5]{\forall a\in X\,(r_e(a)=g)}{\rA}
  \proofstepwidestar[4]{X\subseteq A\land X\neq\varnothing}{%
      \FormulaRefAuto{NonemptyPowersetMembership}[thm]{4}}
  \proofstepwidestarbreakkeep[4]{X\neq\varnothing}{%
    \rAEb{6}}
  \proofstepwidestarbreakkeep[4]{\exists a\,(a\in X)}{%
    \FormulaRefAuto{A\neq\varnothing\eqvdash\exists x\,(x\in A)}{7}}
  \proofstepwidestarbreakkeep[9]{a\in X}{\rA}
  \proofstepwidestarbreakkeep[4,9]{a\in A}{\FormulaRefAuto{PowerSemigroupCarrierElement}[thm]{4,9}}
  \proofstepwidestar[1,2,3,4,9]{e\in A\land(e\star a\in A^2\land a\star e\in A^2)}{%
      \FormulaRefAuto{MonoidSquareIdentityProductCarrier}[thm]{1,2,3,10}}
  \proofstepwidestarbreakkeep[1,2,3,4,9]{e\in A}{%
    \rAEa{11}}
  \proofstepwidestarbreakkeep[1,2,3,4,9]{\{e\}\in\PowNE A}{%
    \FormulaRefAuto{MonoidPowerSingletonCarrier}[thm]{12}}
  \proofstepwidestar[1,2,3,4,9]{e\in\{e\}}{%
      \FormulaRefAuto{a\in\{a\}}}
  \proofstepwidestarbreakkeep[1,2,3,4,9]{e\star a\in\{e\}\star_{\mathcal P}X}{%
    \FormulaRefAuto{PowerSemigroupProductContains}[thm]{1,13,4,14,9}}
  \proofstepwidestarbreakkeep[1,2,3,4,9]{r_e(a)=e\star a}{%
    \FormulaRefAuto{MonoidSquareInflationDef}[def]{1,2,3,10}}
  \proofstepwidestarbreakkeep[5,9]{r_e(a)=g}{\rRE{9,\rUE{5}}}
  \proofstepwidestarbreakkeep[1,2,3,4,5,9]{g\in\{e\}\star_{\mathcal P}X}{\rIE{16,17,15}}
  \proofstepwidestarbreakkeep[1,2,3,4,5]{g\in\{e\}\star_{\mathcal P}X}{\rEE{8,9,18}}
  \proofstepwidestarbreakkeep[20]{z\in\{e\}\star_{\mathcal P}X}{\rA}
  \proofstepwidestarbreakkeep[1,2,3,4,5,20]{z=g}{%
    \FormulaRefAuto{MonoidSquarePowerFiberProductElement}[thm]{1,2,3,4,5,20}}
  \proofstepwidestarbreakkeep[22]{z=g}{\rA}
  \proofstepwidestarbreakkeep[1,2,3,4,5,22]{z\in\{e\}\star_{\mathcal P}X}{\rIE{22,19}}
  \proofstepwidestarbreakkeep[1,2,3,4,5]{\forall z\,(z\in\{e\}\star_{\mathcal P}X\leftrightarrow z=g)}{%
    \rUI{\rLRI{\rRI{20,21},\rRI{22,23}}}}
  \proofstepwidestar[1,2,3,4,5]{z\in\{g\}\leftrightarrow z=g}{%
      \FormulaRefAuto{x\in\{a\}\eqvdash x=a}}
  \proofstepwidestarbreakkeep[1,2,3,4,5]{\{e\}\star_{\mathcal P}X=\{g\}}{%
    \FormulaRefAuto{\forall x\,(x\in A\leftrightarrow x\in B)\eqvdash A=B}{%
      24,25}}
\end{tabproofwide}

\FormulaThmDeltaK[Kriterium für ein einelementiges Rückzugsbild]{%
  \begin{aligned}[t]
    &\SemigroupAlg{A}{\star}\dsep A\neq\varnothing\dsep
      \MonoidAlg{A^2}{\star}{e}\dsep X\in\PowNE A\vdash{}\\[-2pt]
    &\{e\}\star_{\mathcal P}X=\{g\}
      \leftrightarrow\forall a\in X\,(r_e(a)=g)
  \end{aligned}
}{MonoidSquarePowerFiberCriterion}{%
  \DeltaRow{Mengen}{A\dsep e\dsep g\dsep X\dsep a}
  \DeltaRow{Funktionen}{\star\dsep\star_{\mathcal P}\dsep r_e}
}
\begin{tabproofwide}
  \proofstepwidestarbreakkeep[1]{\SemigroupAlg{A}{\star}}{\rA}
  \proofstepwidestarbreakkeep[2]{A\neq\varnothing}{\rA}
  \proofstepwidestarbreakkeep[3]{\MonoidAlg{A^2}{\star}{e}}{\rA}
  \proofstepwidestarbreakkeep[4]{X\in\PowNE A}{\rA}
  \proofstepwidestarbreakkeep[5]{\{e\}\star_{\mathcal P}X=\{g\}}{\rA}
  \proofstepwidestarbreakkeep[6]{a\in X}{\rA}
  \proofstepwidestarbreakkeep[1,2,3,4,5,6]{r_e(a)=g}{%
    \FormulaRefAuto{MonoidSquarePowerFiberElementForward}[thm]{1,2,3,4,5,6}}
  \proofstepwidestarbreakkeep[1,2,3,4,5]{\forall a\in X\,(r_e(a)=g)}{\rUI{\rRI{6,7}}}
  \proofstepwidestarbreakkeep[9]{\forall a\in X\,(r_e(a)=g)}{\rA}
  \proofstepwidestarbreakkeep[1,2,3,4,9]{\{e\}\star_{\mathcal P}X=\{g\}}{%
    \FormulaRefAuto{MonoidSquarePowerFiberBackward}[thm]{1,2,3,4,9}}
  \proofstepwidestarbreakkeep[1,2,3,4]{\{e\}\star_{\mathcal P}X=\{g\}\leftrightarrow\forall a\in X\,(r_e(a)=g)}{%
    \rLRI{\rRI{5,8},\rRI{9,10}}}
\end{tabproofwide}

\par\noindent\begin{minipage}{\linewidth}
\FormulaThmDeltaK[Potenzfasern über dem monoidalen Quadrat]{%
  \begin{aligned}[t]
    &\SemigroupAlg{A}{\star}\dsep A\neq\varnothing\dsep
      \MonoidAlg{A^2}{\star}{e}\dsep g\in A^2\vdash{}\\[-2pt]
    &\bigl\{X\in\PowNE A\bigm|\{e\}\star_{\mathcal P}X=\{g\}\bigr\}\\[-2pt]
    &\qquad=\PowNE{\{a\in A\mid r_e(a)=g\}}
  \end{aligned}
}{MonoidSquarePowerFiber}{%
  \DeltaRow{Mengen}{A\dsep e\dsep g\dsep X\dsep a}
  \DeltaRow{Funktionen}{\star\dsep\star_{\mathcal P}\dsep r_e}
}
\end{minipage}\par
\begin{tabproofwide}
  \proofstepwidestarbreakkeep[1]{\SemigroupAlg{A}{\star}}{\rA}
  \proofstepwidestarbreakkeep[2]{A\neq\varnothing}{\rA}
  \proofstepwidestarbreakkeep[3]{\MonoidAlg{A^2}{\star}{e}}{\rA}
  \proofstepwidestarbreakkeep[4]{g\in A^2}{\rA}
  \proofstepwidestarbreakkeep{\begin{aligned}[t]
      &X\in\{Y\in\PowNE A\mid\{e\}\star_{\mathcal P}Y=\{g\}\}\\[-2pt]
      &\quad\leftrightarrow X\in\PowNE A\land\{e\}\star_{\mathcal P}X=\{g\}
    \end{aligned}}{%
      \FormulaRefAuto{x\in\{u\in A\mid P(u)\}\eqvdash x\in A\land P(x)}}
  \proofstepwidestarbreakkeep[1,2,3]{\leftrightarrow
    X\in\PowNE A\land\forall a\in X\,(r_e(a)=g)}{%
      \FormulaRefAuto{MonoidSquarePowerFiberCriterion}[thm]{1,2,3}}
  \proofstepwidestarbreakkeep[1,2,3]{\leftrightarrow
    X\neq\varnothing\land X\subseteq\{a\in A\mid r_e(a)=g\}}{%
      \FormulaRefAuto{NonemptyPowersetMembership}[thm],\
      \FormulaRefAuto{A\subseteq B\coloneq\forall x\,(x\in A\rightarrow x\in B)}[def],\
      \FormulaRefAuto{x\in\{u\in A\mid P(u)\}\eqvdash x\in A\land P(x)}}
  \proofstepwidestarbreakkeep[1,2,3]{\leftrightarrow
    X\in\PowNE{\{a\in A\mid r_e(a)=g\}}}{%
      \FormulaRefAuto{NonemptyPowersetMembership}[thm]}
  \proofstepwidestarbreakkeep[1,2,3]{%
    \begin{aligned}[t]
      &\bigl\{X\in\PowNE A\bigm|\{e\}\star_{\mathcal P}X=\{g\}\bigr\}\\[-2pt]
      &\qquad=\PowNE{\{a\in A\mid r_e(a)=g\}}
    \end{aligned}}{%
      \FormulaRefAuto{\forall x\,(x\in A\leftrightarrow x\in B)\eqvdash A=B}{\rUI{\rChain{5,6,7,8}}}}
\end{tabproofwide}

\ProofWideReasonsAtRight

\section{Spezielle Definitionen}

\subsection{Endliche Monoide}

\FormulaDefDeltaK[Begriff des endlichen Monoids]{%
  \FiniteMonoidAlg{A}{\star}{e}%
}{FiniteMonoidDef}{%
  \DeltaRow{Mengen}{A\dsep e}
  \DeltaRow{\textbf{Axiome}}{}
  \DeltaPrem{Monoid}{\MonoidAlg{A}{\star}{e}}
  \DeltaPrem{Endlichkeit}{\Finite{A}}
  \DeltaRow{\textbf{Neues Symbol}}{}
  \DeltaPrem{Endliches Monoid}{\FiniteMonoidAlg{A}{\star}{e}}
}

\FormulaAxiomDeltaK[Zugriff auf die Monoidstruktur]{%
  \FiniteMonoidAlg{A}{\star}{e}\vdash\MonoidAlg{A}{\star}{e}%
}{FiniteMonoidBaseAxiom}{%
  \DeltaRow{Mengen}{A\dsep e}
}

\FormulaAxiomDeltaK[Zugriff auf die Endlichkeit]{%
  \FiniteMonoidAlg{A}{\star}{e}\vdash\Finite{A}%
}{FiniteMonoidFinitenessAxiom}{%
  \DeltaRow{Mengen}{A\dsep e}
}

\subsection{Kommutative Monoide}

\FormulaDefDeltaK[Begriff des kommutativen Monoids]{%
  \CommutativeMonoidAlg{A}{\star}{e}%
}{CommutativeMonoidDef}{%
  \DeltaRow{Mengen}{A\dsep e\dsep x\dsep y}
  \DeltaRow{\textbf{Axiome}}{}
  \DeltaPrem{Monoid}{\MonoidAlg{A}{\star}{e}}
  \DeltaRow{Kommutativität}{x,y\in A\vdash x\star y=y\star x}
  \DeltaRow{\textbf{Neues Symbol}}{}
  \DeltaPrem{Kommutatives Monoid}{%
    \CommutativeMonoidAlg{A}{\star}{e}}
}

\FormulaAxiomDeltaK[Zugriff auf das Monoid]{%
  \CommutativeMonoidAlg{A}{\star}{e}
  \vdash\MonoidAlg{A}{\star}{e}%
}{CommutativeMonoidBaseMonoidAxiom}{%
  \DeltaRow{Mengen}{A\dsep e}
}

\FormulaAxiomDeltaK[Zugriff auf die Kommutativität eines Monoids]{%
  \CommutativeMonoidAlg{A}{\star}{e}\dsep x,y\in A
  \vdash x\star y=y\star x%
}{CommutativeMonoidCommutativityAxiom}{%
  \DeltaRow{Mengen}{A\dsep e\dsep x\dsep y}
}

\FormulaThmDeltaK[Kommutative Monoide sind kommutative Halbgruppen]{%
  \CommutativeMonoidAlg{A}{\star}{e}
  \vdash\CommutativeSemigroupAlg{A}{\star}%
}{CommutativeMonoidIsCommutativeSemigroup}{%
  \DeltaRow{Mengen}{A\dsep e\dsep x\dsep y}
}
\begin{tabproof}
  \proofstep{1}{\CommutativeMonoidAlg{A}{\star}{e}}{\rA}
  \proofstep{1}{\MonoidAlg{A}{\star}{e}}{%
    \FormulaRefAuto{CommutativeMonoidBaseMonoidAxiom}[ax]{1}}
  \proofstep{1}{\SemigroupAlg{A}{\star}}{%
    \FormulaRefAuto{MonoidBaseSemigroupAxiom}[ax]{2}}
  \proofstep{4}{x\in A}{\rA}
  \proofstep{5}{y\in A}{\rA}
  \proofstep{1,4,5}{x\star y=y\star x}{%
    \FormulaRefAuto{CommutativeMonoidCommutativityAxiom}[ax]{1,4,5}}
  \proofstep{1}{\CommutativeSemigroupAlg{A}{\star}}{%
    \FormulaRefAuto{CommutativeSemigroupDef}[def]{3,6}}
\end{tabproof}

\section{Einheiten, Teilbarkeit, Irreduzibilität und Primelemente}

\FormulaDefDeltaK[Einheit eines Monoids]{%
  \MonoidUnit{A,\star,e}{u}%
}{MonoidUnitDef}{%
  \DeltaRow{Mengen}{A\dsep e\dsep u\dsep v}
  \DeltaRow{Funktionen}{\star}
  \DeltaPrem{Monoid}{\MonoidAlg{A}{\star}{e}}
  \DeltaRow{Träger}{u\in A}
  \DeltaRow{Inverses Element}{%
    \exists v\in A\,(u\star v=e\land v\star u=e)}
  \DeltaRow{\textbf{Neues Symbol}}{}
  \DeltaPrem{Einheit}{\MonoidUnit{A,\star,e}{u}}
}

\FormulaThmDeltaK[Einheiten einer Potenzhalbgruppe]{%
  \begin{aligned}[t]
    &\MonoidAlg{A}{\star}{e}\dsep X\in\PowNE A\vdash{}\\[-2pt]
    &\MonoidUnit{\PowNE A,\star_{\mathcal P},\{e\}}{X}
      \leftrightarrow
      \exists u\in A\,\bigl(
        X=\{u\}\land\MonoidUnit{A,\star,e}{u}
      \bigr)
  \end{aligned}%
}{PowerMonoidUnitCharacterization}{%
  \DeltaRow{Mengen}{A\dsep e\dsep X\dsep Y\dsep u\dsep v\dsep
    x\dsep x'\dsep y\dsep y'}
  \DeltaRow{Funktionen}{\star\dsep\star_{\mathcal P}}
}
\begin{proof}
\par\medskip\noindent\textbf{(i) Einheiten der Potenzhalbgruppe sind Einermengen von Einheiten.}\par
Sei zunächst \(X\) eine Einheit.  Nach
\FormulaRefAuto{MonoidUnitDef}[def] gibt es ein \(Y\in\PowNE A\) mit
\[
  X\star_{\mathcal P}Y=\{e\}
  =Y\star_{\mathcal P}X.
\]
Wähle \(x,x'\in X\) und \(y\in Y\).  Dann gilt
\(x\star y=x'\star y=e\) sowie
\(y\star x=y\star x'=e\).  Mit den Monoidgesetzen folgt
\[
  x'=x'\star e
     =x'\star(y\star x)
     =(x'\star y)\star x
     =e\star x=x.
\]
Also ist \(X=\{u\}\) für ein \(u\in A\).  Für \(y,y'\in Y\) und
\(x\in X\) gilt ebenso
\[
  y'=e\star y'
     =(y\star x)\star y'
     =y\star(x\star y')
     =y\star e=y.
\]
Somit ist \(Y=\{v\}\) für ein \(v\in A\).  Die Produktgleichungen ergeben
\(u\star v=e=v\star u\), also ist \(u\) eine Einheit von \(A\).

\par\medskip\noindent\textbf{(ii) Einermengen von Einheiten sind Einheiten.}\par
Umgekehrt sei \(X=\{u\}\) und \(u\) eine Einheit von \(A\).  Für ein
inverses Element \(v\in A\) gilt
\[
  \{u\}\star_{\mathcal P}\{v\}=\{e\}
  =\{v\}\star_{\mathcal P}\{u\}.
\]
Nach \FormulaRefAuto{PowerSemigroupMonoid}[thm] ist
\((\PowNE A,\star_{\mathcal P},\{e\})\) ein Monoid; die beiden
Produktgleichungen erfüllen daher
\FormulaRefAuto{MonoidUnitDef}[def].  Somit ist \(X\) eine Einheit der
Potenzhalbgruppe.
\end{proof}

\FormulaThmDeltaK[Das neutrale Element ist eine Einheit]{%
  \MonoidAlg{A}{\star}{e}
  \vdash\MonoidUnit{A,\star,e}{e}%
}{MonoidIdentityIsUnit}{%
  \DeltaRow{Mengen}{A\dsep e\dsep v}
  \DeltaRow{Funktionen}{\star}
}
\begin{tabproof}
  \proofstep{1}{\MonoidAlg{A}{\star}{e}}{\rA}
  \proofstep{1}{e\in A}{%
    \FormulaRefAuto{MonoidIdentityCarrierAxiom}[ax]{1}}
  \proofstep{1}{e\star e=e}{%
    \FormulaRefAuto{MonoidLeftIdentityAxiom}[ax]{1,2}}
  \proofstep{1}{%
    \exists v\in A\,(e\star v=e\land v\star e=e)}{%
    \rEI{2,\rAIReuse{2}{3}}}
  \proofstep{1}{\MonoidUnit{A,\star,e}{e}}{%
    \FormulaRefAuto{MonoidUnitDef}[def]{1,2,4}}
\end{tabproof}

\paragraph{Beispiel.}
Das neutrale Element muss nicht die einzige Einheit sein. Im additiven Monoid
\((\mathbb Z,+,\IntZero)\) ist etwa \(\IntOne\) eine vom
neutralen Element verschiedene Einheit. Ihr inverses Element ist
\(-\IntOne\), denn nach \FormulaRefAuto{IntegerAddInverseAxiom}[ax] gilt
\[
  \IntOne+(-\IntOne)=\IntZero
  \qquad\text{und}\qquad
  (-\IntOne)+\IntOne=\IntZero.
\]
Dabei ist \(\IntOne\neq\IntZero\). Tatsächlich ist im additiven
Monoid der ganzen Zahlen jede ganze Zahl eine Einheit.

\FormulaDefDeltaK[Teilbarkeit in einem kommutativen Monoid]{%
  \MonoidDivides{A,\star,e}{a}{b}%
}{CommutativeMonoidDividesDef}{%
  \DeltaRow{Mengen}{A\dsep a\dsep b\dsep c\dsep e}
  \DeltaRow{Funktionen}{\star}
  \DeltaPrem{Kommutatives Monoid}{%
    \CommutativeMonoidAlg{A}{\star}{e}}
  \DeltaRow{Elemente}{a,b\in A}
  \DeltaRow{Quotient}{\exists c\in A\,(b=a\star c)}
  \DeltaRow{\textbf{Neues Symbol}}{}
  \DeltaPrem{Teilbarkeit}{\MonoidDivides{A,\star,e}{a}{b}}
}

\FormulaThmDeltaK[Reflexivität der Teilbarkeit]{%
  \CommutativeMonoidAlg{A}{\star}{e}\dsep a\in A
  \vdash\MonoidDivides{A,\star,e}{a}{a}%
}{CommutativeMonoidDividesReflexive}{%
  \DeltaRow{Mengen}{A\dsep a\dsep c\dsep e}
  \DeltaRow{Funktionen}{\star}
}
\begin{tabproof}
  \proofstep{1}{\CommutativeMonoidAlg{A}{\star}{e}}{\rA}
  \proofstep{2}{a\in A}{\rA}
  \proofstep{1}{\MonoidAlg{A}{\star}{e}}{%
    \FormulaRefAuto{CommutativeMonoidBaseMonoidAxiom}[ax]{1}}
  \proofstep{1}{e\in A}{%
    \FormulaRefAuto{MonoidIdentityCarrierAxiom}[ax]{3}}
  \proofstep{1,2}{a\star e=a}{%
    \FormulaRefAuto{MonoidRightIdentityAxiom}[ax]{3,2}}
  \proofstep{1,2}{a=a\star e}{%
    \FormulaRefAuto{a=b\vdash b=a}{5}}
  \proofstep{1,2}{\exists c\in A\,(a=a\star c)}{%
    \rEI{4,6}}
  \proofstep{1,2}{\MonoidDivides{A,\star,e}{a}{a}}{%
    \FormulaRefAuto{CommutativeMonoidDividesDef}[def]{1,2,2,7}}
\end{tabproof}

\FormulaThmDeltaK[Transitivität der Teilbarkeit]{%
  \begin{aligned}[t]
    &\CommutativeMonoidAlg{A}{\star}{e}\dsep a,b,c\in A\dsep
      \MonoidDivides{A,\star,e}{a}{b}\dsep
      \MonoidDivides{A,\star,e}{b}{c}\vdash{}\\[-2pt]
    &\MonoidDivides{A,\star,e}{a}{c}
  \end{aligned}
}{CommutativeMonoidDividesTransitive}{%
  \DeltaRow{Mengen}{A\dsep a\dsep b\dsep c\dsep e\dsep r\dsep s}
  \DeltaRow{Funktionen}{\star}
}
\begin{tabproofsplitwide}
  \proofpartwideindR[Feste Teilbarkeitszeugen verknüpfen]{%
    \begin{aligned}[t]
      &\CommutativeMonoidAlg{A}{\star}{e}\dsep
        a,b,c,r,s\in A\dsep b=a\star r\dsep c=b\star s\vdash{}\\[-2pt]
      &\MonoidDivides{A,\star,e}{a}{c}
    \end{aligned}
  }{%
    \begin{aligned}[t]
      &\CommutativeMonoidAlg{A}{\star}{e}\dsep
        a,b,c,r,s\in A\dsep b=a\star r\dsep c=b\star s\vdash{}\\[-2pt]
      &\MonoidDivides{A,\star,e}{a}{c}
    \end{aligned}
  }[CommutativeMonoidDividesTransitiveFixedWitnesses]
    \proofstepwidestar[1]{\CommutativeMonoidAlg{A}{\star}{e}}{\rA}
    \proofstepwidestar[2]{a\in A}{\rA}
    \proofstepwidestar[3]{b\in A}{\rA}
    \proofstepwidestar[4]{c\in A}{\rA}
    \proofstepwidestar[5]{r\in A}{\rA}
    \proofstepwidestar[6]{s\in A}{\rA}
    \proofstepwidestar[7]{b=a\star r}{\rA}
    \proofstepwidestar[8]{c=b\star s}{\rA}
    \proofstepwidestar[1]{\MonoidAlg{A}{\star}{e}}{%
      \FormulaRefAuto{CommutativeMonoidBaseMonoidAxiom}[ax]{1}}
    \proofstepwidestar[1]{\SemigroupAlg{A}{\star}}{%
      \FormulaRefAuto{MonoidBaseSemigroupAxiom}[ax]{9}}
    \proofstepwidestar[1,5,6]{r\star s\in A}{%
      \FormulaRefAuto{SemigroupClosureAxiom}[ax]{10,5,6}}
    \proofstepwide[8]{c}{=}{b\star s}{\rA}
    \proofstepwide[6,7]{}{=}{(a\star r)\star s}{\rIE{7}}
    \proofstepwide[1,2,5,6]{}{=}{a\star(r\star s)}{%
      \FormulaRefAuto{SemigroupAssociativityAxiom}[ax]{10,2,5,6}}
    \proofstepwidestar[1,2,5,6,7,8]{c=a\star(r\star s)}{%
      \rChain{12,14}}
    \proofstepwidestar[1,2,5,6,7,8]{%
      \exists r\in A\,(c=a\star r)}{\rEI{11,15}}
    \proofstepwidestar[1,2,4,5,6,7,8]{%
      \MonoidDivides{A,\star,e}{a}{c}}{%
      \FormulaRefAuto{CommutativeMonoidDividesDef}[def]{1,2,4,16}}
  \closeproofpartwideindR

  \proofpartwide{Elimination der Teilbarkeitszeugen}
    \proofstepwidestar[1]{\CommutativeMonoidAlg{A}{\star}{e}}{\rA}
    \proofstepwidestar[2]{a\in A}{\rA}
    \proofstepwidestar[3]{b\in A}{\rA}
    \proofstepwidestar[4]{c\in A}{\rA}
    \proofstepwidestar[5]{\MonoidDivides{A,\star,e}{a}{b}}{\rA}
    \proofstepwidestar[6]{\MonoidDivides{A,\star,e}{b}{c}}{\rA}
    \proofstepwidestar[5]{\exists r\in A\,(b=a\star r)}{%
      \FormulaRefAuto{CommutativeMonoidDividesDef}[def]{5}}
    \proofstepwidestar[6]{\exists s\in A\,(c=b\star s)}{%
      \FormulaRefAuto{CommutativeMonoidDividesDef}[def]{6}}
    \proofstepwidestar[9]{r\in A}{\rA}
    \proofstepwidestar[10]{b=a\star r}{\rA}
    \proofstepwidestar[11]{s\in A}{\rA}
    \proofstepwidestar[12]{c=b\star s}{\rA}
    \proofstepwidestar[1,2,3,4,9,10,11,12]{%
      \MonoidDivides{A,\star,e}{a}{c}}{%
      \FormulaRefAuto{CommutativeMonoidDividesTransitiveFixedWitnesses}[thm]{%
        1,2,3,4,9,11,10,12}}
    \proofstepwidestar[1,2,3,4,5,6]{%
      \MonoidDivides{A,\star,e}{a}{c}}{%
      \rEE{7,9,10,\rEE{8,11,12,13}}}
  \closeproofpartwide
\end{tabproofsplitwide}

\FormulaThmDeltaK[Einheiten sind genau die Teiler des neutralen Elements]{%
  \begin{aligned}[t]
    &\CommutativeMonoidAlg{A}{\star}{e}\dsep u\in A\vdash{}\\[-2pt]
    &\MonoidUnit{A,\star,e}{u}
      \leftrightarrow\MonoidDivides{A,\star,e}{u}{e}
  \end{aligned}
}{CommutativeMonoidUnitIffDividesIdentity}{%
  \DeltaRow{Mengen}{A\dsep e\dsep u\dsep v\dsep w}
  \DeltaRow{Funktionen}{\star}
}
\begin{tabproofsplitwide}
  \proofpartwideindR[Eine Einheit teilt das neutrale Element]{%
    \begin{aligned}[t]
      &\CommutativeMonoidAlg{A}{\star}{e}\dsep u\in A\dsep
        \MonoidUnit{A,\star,e}{u}\vdash{}\\[-2pt]
      &\MonoidDivides{A,\star,e}{u}{e}
    \end{aligned}
  }{%
    \begin{aligned}[t]
      &\CommutativeMonoidAlg{A}{\star}{e}\dsep u\in A\dsep
        \MonoidUnit{A,\star,e}{u}\vdash{}\\[-2pt]
      &\MonoidDivides{A,\star,e}{u}{e}
    \end{aligned}
  }[CommutativeMonoidUnitImpliesDividesIdentity]
    \proofstepwidestar[1]{\CommutativeMonoidAlg{A}{\star}{e}}{\rA}
    \proofstepwidestar[2]{u\in A}{\rA}
    \proofstepwidestar[3]{\MonoidUnit{A,\star,e}{u}}{\rA}
    \proofstepwidestar[1]{\MonoidAlg{A}{\star}{e}}{%
      \FormulaRefAuto{CommutativeMonoidBaseMonoidAxiom}[ax]{1}}
    \proofstepwidestar[1]{e\in A}{%
      \FormulaRefAuto{MonoidIdentityCarrierAxiom}[ax]{4}}
    \proofstepwidestar[3]{%
      \exists v\in A\,(u\star v=e\land v\star u=e)}{%
      \FormulaRefAuto{MonoidUnitDef}[def]{3}}
    \proofstepwidestar[7]{v\in A}{\rA}
    \proofstepwidestar[8]{u\star v=e\land v\star u=e}{\rA}
    \proofstepwidestar[8]{u\star v=e}{\rAEa{8}}
    \proofstepwidestar[8]{e=u\star v}{%
      \FormulaRefAuto{a=b\vdash b=a}{9}}
    \proofstepwidestar[7,8]{\exists v\in A\,(e=u\star v)}{%
      \rEI{7,10}}
    \proofstepwidestar[1,2,7,8]{%
      \MonoidDivides{A,\star,e}{u}{e}}{%
      \FormulaRefAuto{CommutativeMonoidDividesDef}[def]{1,2,5,11}}
    \proofstepwidestar[1,2,3]{\MonoidDivides{A,\star,e}{u}{e}}{%
      \rEE{6,7,8,12}}
  \closeproofpartwideindR

  \proofpartwideindR[Ein Teiler des neutralen Elements ist eine Einheit]{%
    \begin{aligned}[t]
      &\CommutativeMonoidAlg{A}{\star}{e}\dsep u\in A\dsep
        \MonoidDivides{A,\star,e}{u}{e}\vdash{}\\[-2pt]
      &\MonoidUnit{A,\star,e}{u}
    \end{aligned}
  }{%
    \begin{aligned}[t]
      &\CommutativeMonoidAlg{A}{\star}{e}\dsep u\in A\dsep
        \MonoidDivides{A,\star,e}{u}{e}\vdash{}\\[-2pt]
      &\MonoidUnit{A,\star,e}{u}
    \end{aligned}
  }[CommutativeMonoidDividesIdentityImpliesUnit]
    \proofstepwidestar[1]{\CommutativeMonoidAlg{A}{\star}{e}}{\rA}
    \proofstepwidestar[2]{u\in A}{\rA}
    \proofstepwidestar[3]{\MonoidDivides{A,\star,e}{u}{e}}{\rA}
    \proofstepwidestar[1]{\MonoidAlg{A}{\star}{e}}{%
      \FormulaRefAuto{CommutativeMonoidBaseMonoidAxiom}[ax]{1}}
    \proofstepwidestar[1]{e\in A}{%
      \FormulaRefAuto{MonoidIdentityCarrierAxiom}[ax]{4}}
    \proofstepwidestar[3]{\exists w\in A\,(e=u\star w)}{%
      \FormulaRefAuto{CommutativeMonoidDividesDef}[def]{3}}
    \proofstepwidestar[7]{w\in A}{\rA}
    \proofstepwidestar[8]{e=u\star w}{\rA}
    \proofstepwidestar[8]{u\star w=e}{%
      \FormulaRefAuto{a=b\vdash b=a}{8}}
    \proofstepwidestar[1,2,7]{w\star u=u\star w}{%
      \FormulaRefAuto{CommutativeMonoidCommutativityAxiom}[ax]{1,7,2}}
    \proofstepwidestar[1,2,7,8]{w\star u=e}{\rChain{10,9}}
    \proofstepwidestar[1,2,7,8]{%
      \exists w\in A\,(u\star w=e\land w\star u=e)}{%
      \rEI{7,\rAI{9,11}}}
    \proofstepwidestar[1,2,7,8]{\MonoidUnit{A,\star,e}{u}}{%
      \FormulaRefAuto{MonoidUnitDef}[def]{4,2,12}}
    \proofstepwidestar[1,2,3]{\MonoidUnit{A,\star,e}{u}}{%
      \rEE{6,7,8,13}}
  \closeproofpartwideindR

  \proofpartwide{Zusammenführung der Richtungen}
    \proofstepwidestar[1]{\CommutativeMonoidAlg{A}{\star}{e}}{\rA}
    \proofstepwidestar[2]{u\in A}{\rA}
    \proofstepwidestar[3]{\MonoidUnit{A,\star,e}{u}}{\rA}
    \proofstepwidestar[1,2,3]{\MonoidDivides{A,\star,e}{u}{e}}{%
      \FormulaRefAuto{CommutativeMonoidUnitImpliesDividesIdentity}[thm]{1,2,3}}
    \proofstepwidestar[1,2]{%
      \MonoidUnit{A,\star,e}{u}\rightarrow
      \MonoidDivides{A,\star,e}{u}{e}}{\rRI{3,4}}
    \proofstepwidestar[6]{\MonoidDivides{A,\star,e}{u}{e}}{\rA}
    \proofstepwidestar[1,2,6]{\MonoidUnit{A,\star,e}{u}}{%
      \FormulaRefAuto{CommutativeMonoidDividesIdentityImpliesUnit}[thm]{1,2,6}}
    \proofstepwidestar[1,2]{%
      \MonoidDivides{A,\star,e}{u}{e}\rightarrow
      \MonoidUnit{A,\star,e}{u}}{\rRI{6,7}}
    \proofstepwidestar[1,2]{%
      \MonoidUnit{A,\star,e}{u}
      \leftrightarrow\MonoidDivides{A,\star,e}{u}{e}}{%
      \rLRI{5,8}}
  \closeproofpartwide
\end{tabproofsplitwide}

\FormulaDefDeltaK[Irreduzibles Element eines kommutativen Monoids]{%
  \MonoidIrreducible{A,\star,e}{p}%
}{CommutativeMonoidIrreducibleDef}{%
  \DeltaRow{Mengen}{A\dsep a\dsep b\dsep e\dsep p}
  \DeltaRow{Funktionen}{\star}
  \DeltaPrem{Kommutatives Monoid}{%
    \CommutativeMonoidAlg{A}{\star}{e}}
  \DeltaRow{Element}{p\in A}
  \DeltaRow{Keine Einheit}{\neg\MonoidUnit{A,\star,e}{p}}
  \DeltaRow{Irreduzibilität}{%
    \begin{aligned}[t]
      &a,b\in A\dsep p=a\star b\vdash{}\\[-2pt]
      &\MonoidUnit{A,\star,e}{a}\lor
       \MonoidUnit{A,\star,e}{b}
    \end{aligned}}
  \DeltaRow{\textbf{Neues Symbol}}{}
  \DeltaPrem{Irreduzibles Element}{%
    \MonoidIrreducible{A,\star,e}{p}}
}

\FormulaDefDeltaK[Primelement eines kommutativen Monoids]{%
  \MonoidPrime{A,\star,e}{p}%
}{CommutativeMonoidPrimeDef}{%
  \DeltaRow{Mengen}{A\dsep a\dsep b\dsep e\dsep p}
  \DeltaRow{Funktionen}{\star}
  \DeltaPrem{Kommutatives Monoid}{%
    \CommutativeMonoidAlg{A}{\star}{e}}
  \DeltaRow{Element}{p\in A}
  \DeltaRow{Keine Einheit}{\neg\MonoidUnit{A,\star,e}{p}}
  \DeltaRow{Primeigenschaft}{%
    \begin{aligned}[t]
      &a,b\in A\dsep
        \MonoidDivides{A,\star,e}{p}{a\star b}\vdash{}\\[-2pt]
      &\MonoidDivides{A,\star,e}{p}{a}\lor
       \MonoidDivides{A,\star,e}{p}{b}
    \end{aligned}}
  \DeltaRow{\textbf{Neues Symbol}}{}
  \DeltaPrem{Primelement}{\MonoidPrime{A,\star,e}{p}}
}

\FormulaThmDeltaK[Ein geteilter Faktor erzwingt eine Einheit]{%
  \begin{aligned}[t]
    &\CommutativeMonoidAlg{A}{\star}{e}\dsep
      \mathsf{K\ddot{u}rz}\!\left(A,\star\right)\dsep
      p,a,b\in A\dsep p=a\star b\dsep{}\\[-2pt]
    &\MonoidDivides{A,\star,e}{p}{a}
      \vdash\MonoidUnit{A,\star,e}{b}
  \end{aligned}
}{CommutativeCancellativeMonoidFactorDivisorUnit}{%
  \DeltaRow{Mengen}{A\dsep a\dsep b\dsep c\dsep e\dsep p}
  \DeltaRow{Funktionen}{\star}
}
\begin{tabproofsplitwide}
  \proofpartwideindR[Herstellung der Kürzungsgleichung]{%
    \begin{aligned}[t]
      &\CommutativeMonoidAlg{A}{\star}{e}\dsep
        p,a,b,c\in A\dsep p=a\star b\dsep a=p\star c\vdash{}\\[-2pt]
      &p\star e=p\star(c\star b)
    \end{aligned}
  }{%
    \begin{aligned}[t]
      &\CommutativeMonoidAlg{A}{\star}{e}\dsep
        p,a,b,c\in A\dsep p=a\star b\dsep a=p\star c\vdash{}\\[-2pt]
      &p\star e=p\star(c\star b)
    \end{aligned}
  }[CommutativeMonoidFactorDivisorCancellationEquation]
    \proofstepwidestar[1]{\CommutativeMonoidAlg{A}{\star}{e}}{\rA}
    \proofstepwidestar[2]{p\in A}{\rA}
    \proofstepwidestar[3]{a\in A}{\rA}
    \proofstepwidestar[4]{b\in A}{\rA}
    \proofstepwidestar[5]{c\in A}{\rA}
    \proofstepwidestar[6]{p=a\star b}{\rA}
    \proofstepwidestar[7]{a=p\star c}{\rA}
    \proofstepwidestar[1]{\MonoidAlg{A}{\star}{e}}{%
      \FormulaRefAuto{CommutativeMonoidBaseMonoidAxiom}[ax]{1}}
    \proofstepwidestar[1]{\SemigroupAlg{A}{\star}}{%
      \FormulaRefAuto{MonoidBaseSemigroupAxiom}[ax]{8}}
    \proofstepwide[1,2]{p\star e}{=}{p}{%
      \FormulaRefAuto{MonoidRightIdentityAxiom}[ax]{8,2}}
    \proofstepwide[6]{}{=}{a\star b}{\rA}
    \proofstepwide[4,7]{}{=}{(p\star c)\star b}{\rIE{7}}
    \proofstepwide[1,2,4,5]{}{=}{p\star(c\star b)}{%
      \FormulaRefAuto{SemigroupAssociativityAxiom}[ax]{9,2,5,4}}
    \proofstepwidestar[1,2,4,5,6,7]{%
      p\star e=p\star(c\star b)}{\rChain{10,13}}
  \closeproofpartwideindR

  \proofpartwideindR[Kürzung und Gewinnung des Inversen]{%
    \begin{aligned}[t]
      &\CommutativeMonoidAlg{A}{\star}{e}\dsep
        \mathsf{K\ddot{u}rz}\!\left(A,\star\right)\dsep
        p,b,c\in A\dsep{}\\[-2pt]
      &p\star e=p\star(c\star b)
        \vdash\MonoidUnit{A,\star,e}{b}
    \end{aligned}
  }{%
    \begin{aligned}[t]
      &\CommutativeMonoidAlg{A}{\star}{e}\dsep
        \mathsf{K\ddot{u}rz}\!\left(A,\star\right)\dsep
        p,b,c\in A\dsep{}\\[-2pt]
      &p\star e=p\star(c\star b)
        \vdash\MonoidUnit{A,\star,e}{b}
    \end{aligned}
  }[CommutativeCancellativeMonoidCancellationEquationImpliesUnit]
    \proofstepwidestar[1]{\CommutativeMonoidAlg{A}{\star}{e}}{\rA}
    \proofstepwidestar[2]{\mathsf{K\ddot{u}rz}\!\left(A,\star\right)}{\rA}
    \proofstepwidestar[3]{p\in A}{\rA}
    \proofstepwidestar[4]{b\in A}{\rA}
    \proofstepwidestar[5]{c\in A}{\rA}
    \proofstepwidestar[6]{p\star e=p\star(c\star b)}{\rA}
    \proofstepwidestar[1]{\MonoidAlg{A}{\star}{e}}{%
      \FormulaRefAuto{CommutativeMonoidBaseMonoidAxiom}[ax]{1}}
    \proofstepwidestar[1]{\SemigroupAlg{A}{\star}}{%
      \FormulaRefAuto{MonoidBaseSemigroupAxiom}[ax]{7}}
    \proofstepwidestar[1]{e\in A}{%
      \FormulaRefAuto{MonoidIdentityCarrierAxiom}[ax]{7}}
    \proofstepwidestar[2]{\mathsf{LK\ddot{u}rz}\!\left(A,\star\right)}{%
      \FormulaRefAuto{CancellativeLeftAxiom}[ax]{2}}
    \proofstepwidestar[1,4,5]{c\star b\in A}{%
      \FormulaRefAuto{SemigroupClosureAxiom}[ax]{8,5,4}}
    \proofstepwidestar[1,2,3,4,5,6]{e=c\star b}{%
      \FormulaRefAuto{LeftCancellationAxiom}[ax]{10,3,9,11,6}}
    \proofstepwidestar[1,2,3,4,5,6]{c\star b=e}{%
      \FormulaRefAuto{a=b\vdash b=a}{12}}
    \proofstepwide[1,4,5]{b\star c}{=}{c\star b}{%
      \FormulaRefAuto{CommutativeMonoidCommutativityAxiom}[ax]{1,4,5}}
    \proofstepwidestar[1,2,3,4,5,6]{b\star c=e}{%
      \rChain{14,13}}
    \proofstepwidestar[1,2,3,4,5,6]{%
      \exists c\in A\,(b\star c=e\land c\star b=e)}{%
      \rEI{5,\rAI{15,13}}}
    \proofstepwidestar[1,2,3,4,5,6]{\MonoidUnit{A,\star,e}{b}}{%
      \FormulaRefAuto{MonoidUnitDef}[def]{7,4,16}}
  \closeproofpartwideindR

  \proofpartwide{Elimination des Teilbarkeitszeugen}
    \proofstepwidestar[1]{\CommutativeMonoidAlg{A}{\star}{e}}{\rA}
    \proofstepwidestar[2]{\mathsf{K\ddot{u}rz}\!\left(A,\star\right)}{\rA}
    \proofstepwidestar[3]{p\in A}{\rA}
    \proofstepwidestar[4]{a\in A}{\rA}
    \proofstepwidestar[5]{b\in A}{\rA}
    \proofstepwidestar[6]{p=a\star b}{\rA}
    \proofstepwidestar[7]{\MonoidDivides{A,\star,e}{p}{a}}{\rA}
    \proofstepwidestar[7]{\exists c\in A\,(a=p\star c)}{%
      \FormulaRefAuto{CommutativeMonoidDividesDef}[def]{7}}
    \proofstepwidestar[9]{c\in A}{\rA}
    \proofstepwidestar[10]{a=p\star c}{\rA}
    \proofstepwidestar[1,3,4,5,6,9,10]{p\star e=p\star(c\star b)}{%
      \FormulaRefAuto{CommutativeMonoidFactorDivisorCancellationEquation}[thm]{%
        1,3,4,5,9,6,10}}
    \proofstepwidestar[1,2,3,4,5,6,9,10]{\MonoidUnit{A,\star,e}{b}}{%
      \FormulaRefAuto{CommutativeCancellativeMonoidCancellationEquationImpliesUnit}[thm]{%
        1,2,3,5,9,11}}
    \proofstepwidestar[1,2,3,4,5,6,7]{\MonoidUnit{A,\star,e}{b}}{%
      \rEE{8,9,10,12}}
  \closeproofpartwide
\end{tabproofsplitwide}

\FormulaThmDeltaK[Primelemente in kommutativen kürzbaren Monoiden sind irreduzibel]{%
  \begin{aligned}[t]
    &\CommutativeMonoidAlg{A}{\star}{e}\dsep
      \mathsf{K\ddot{u}rz}\!\left(A,\star\right)\dsep
      \MonoidPrime{A,\star,e}{p}\vdash{}\\[-2pt]
    &\MonoidIrreducible{A,\star,e}{p}
  \end{aligned}
}{CommutativeCancellativeMonoidPrimeIsIrreducible}{%
  \DeltaRow{Mengen}{A\dsep a\dsep b\dsep e\dsep p}
  \DeltaRow{Funktionen}{\star}
}
\begin{tabproofsplitwide}
  \proofpartwideindR[Teilerdisjunktion der Faktorisierung]{%
    \begin{aligned}[t]
      &\CommutativeMonoidAlg{A}{\star}{e}\dsep
        \MonoidPrime{A,\star,e}{p}\dsep{}\\[-2pt]
      &a,b\in A\dsep p=a\star b\vdash{}\\[-2pt]
      &\MonoidDivides{A,\star,e}{p}{a}\lor
       \MonoidDivides{A,\star,e}{p}{b}
    \end{aligned}
  }{%
    \begin{aligned}[t]
      &\CommutativeMonoidAlg{A}{\star}{e}\dsep
        \MonoidPrime{A,\star,e}{p}\dsep{}\\[-2pt]
      &a,b\in A\dsep p=a\star b\vdash{}\\[-2pt]
      &\MonoidDivides{A,\star,e}{p}{a}\lor
       \MonoidDivides{A,\star,e}{p}{b}
    \end{aligned}
  }[CommutativeMonoidPrimeFactorDivisorDisjunction]
    \proofstepwidestar[1]{\CommutativeMonoidAlg{A}{\star}{e}}{\rA}
    \proofstepwidestar[2]{\MonoidPrime{A,\star,e}{p}}{\rA}
    \proofstepwidestar[3]{a\in A}{\rA}
    \proofstepwidestar[4]{b\in A}{\rA}
    \proofstepwidestar[5]{p=a\star b}{\rA}
    \proofstepwidestar[2]{p\in A}{%
      \FormulaRefAuto{CommutativeMonoidPrimeDef}[def]{2}}
    \proofstepwidestar[1,2]{\MonoidDivides{A,\star,e}{p}{p}}{%
      \FormulaRefAuto{CommutativeMonoidDividesReflexive}[thm]{1,6}}
    \proofstepwidestar[1,2,5]{%
      \MonoidDivides{A,\star,e}{p}{a\star b}}{\rIE{5}}
    \proofstepwidestar[1,2,3,4,5]{%
      \MonoidDivides{A,\star,e}{p}{a}\lor
      \MonoidDivides{A,\star,e}{p}{b}}{%
      \FormulaRefAuto{CommutativeMonoidPrimeDef}[def]{2,3,4,8}}
  \closeproofpartwideindR

  \proofpartwideindR[Von der Teilerdisjunktion zur Einheitendisjunktion]{%
    \begin{aligned}[t]
      &\CommutativeMonoidAlg{A}{\star}{e}\dsep
        \mathsf{K\ddot{u}rz}\!\left(A,\star\right)\dsep
        \MonoidPrime{A,\star,e}{p}\dsep a,b\in A\dsep{}\\[-2pt]
      &p=a\star b\dsep
        \bigl(\MonoidDivides{A,\star,e}{p}{a}\lor
        \MonoidDivides{A,\star,e}{p}{b}\bigr)\vdash{}\\[-2pt]
      &\MonoidUnit{A,\star,e}{a}\lor\MonoidUnit{A,\star,e}{b}
    \end{aligned}
  }{%
    \begin{aligned}[t]
      &\CommutativeMonoidAlg{A}{\star}{e}\dsep
        \mathsf{K\ddot{u}rz}\!\left(A,\star\right)\dsep
        \MonoidPrime{A,\star,e}{p}\dsep a,b\in A\dsep{}\\[-2pt]
      &p=a\star b\dsep
        \bigl(\MonoidDivides{A,\star,e}{p}{a}\lor
        \MonoidDivides{A,\star,e}{p}{b}\bigr)\vdash{}\\[-2pt]
      &\MonoidUnit{A,\star,e}{a}\lor\MonoidUnit{A,\star,e}{b}
    \end{aligned}
  }[CommutativeCancellativeMonoidPrimeFactorUnitDisjunction]
    \proofstepwidestar[1]{\CommutativeMonoidAlg{A}{\star}{e}}{\rA}
    \proofstepwidestar[2]{\mathsf{K\ddot{u}rz}\!\left(A,\star\right)}{\rA}
    \proofstepwidestar[3]{\MonoidPrime{A,\star,e}{p}}{\rA}
    \proofstepwidestar[4]{a\in A}{\rA}
    \proofstepwidestar[5]{b\in A}{\rA}
    \proofstepwidestar[6]{p=a\star b}{\rA}
    \proofstepwidestar[7]{%
      \MonoidDivides{A,\star,e}{p}{a}\lor
      \MonoidDivides{A,\star,e}{p}{b}}{\rA}
    \proofstepwidestar[3]{p\in A}{%
      \FormulaRefAuto{CommutativeMonoidPrimeDef}[def]{3}}
    \proofstepwidestar[9]{\MonoidDivides{A,\star,e}{p}{a}}{\rA}
    \proofstepwidestar[1,2,3,4,5,6,9]{\MonoidUnit{A,\star,e}{b}}{%
      \FormulaRefAuto{CommutativeCancellativeMonoidFactorDivisorUnit}[thm]{%
        1,2,8,4,5,6,9}}
    \proofstepwidestar[1,2,3,4,5,6,9]{%
      \MonoidUnit{A,\star,e}{a}\lor\MonoidUnit{A,\star,e}{b}}{%
      \rOIb{10}}
    \proofstepwidestar[12]{\MonoidDivides{A,\star,e}{p}{b}}{\rA}
    \proofstepwidestar[1,4,5]{a\star b=b\star a}{%
      \FormulaRefAuto{CommutativeMonoidCommutativityAxiom}[ax]{1,4,5}}
    \proofstepwidestar[1,4,5,6]{p=b\star a}{\rChain{6,13}}
    \proofstepwidestar[1,2,3,4,5,6,12]{\MonoidUnit{A,\star,e}{a}}{%
      \FormulaRefAuto{CommutativeCancellativeMonoidFactorDivisorUnit}[thm]{%
        1,2,8,5,4,14,12}}
    \proofstepwidestar[1,2,3,4,5,6,12]{%
      \MonoidUnit{A,\star,e}{a}\lor\MonoidUnit{A,\star,e}{b}}{%
      \rOIa{15}}
    \proofstepwidestar[1,2,3,4,5,6,7]{%
      \MonoidUnit{A,\star,e}{a}\lor\MonoidUnit{A,\star,e}{b}}{%
      \rOE{7,9,11,12,16}}
  \closeproofpartwideindR

  \proofpartwide{Nachweis der Irreduzibilität}
    \proofstepwidestar[1]{\CommutativeMonoidAlg{A}{\star}{e}}{\rA}
    \proofstepwidestar[2]{\mathsf{K\ddot{u}rz}\!\left(A,\star\right)}{\rA}
    \proofstepwidestar[3]{\MonoidPrime{A,\star,e}{p}}{\rA}
    \proofstepwidestar[3]{p\in A}{%
      \FormulaRefAuto{CommutativeMonoidPrimeDef}[def]{3}}
    \proofstepwidestar[3]{\neg\MonoidUnit{A,\star,e}{p}}{%
      \FormulaRefAuto{CommutativeMonoidPrimeDef}[def]{3}}
    \proofstepwidestar[6]{a\in A}{\rA}
    \proofstepwidestar[7]{b\in A}{\rA}
    \proofstepwidestar[8]{p=a\star b}{\rA}
    \proofstepwidestar[1,3,6,7,8]{%
      \MonoidDivides{A,\star,e}{p}{a}\lor
      \MonoidDivides{A,\star,e}{p}{b}}{%
      \FormulaRefAuto{CommutativeMonoidPrimeFactorDivisorDisjunction}[thm]{%
        1,3,6,7,8}}
    \proofstepwidestar[1,2,3,6,7,8]{%
      \MonoidUnit{A,\star,e}{a}\lor\MonoidUnit{A,\star,e}{b}}{%
      \FormulaRefAuto{CommutativeCancellativeMonoidPrimeFactorUnitDisjunction}[thm]{%
        1,2,3,6,7,8,9}}
    \proofstepwidestar[1,2,3]{\MonoidIrreducible{A,\star,e}{p}}{%
      \FormulaRefAuto{CommutativeMonoidIrreducibleDef}[def]{%
        1,4,5,6,7,8,10}}
  \closeproofpartwide
\end{tabproofsplitwide}

\section{Morphismen von Monoiden}

\FormulaDefDeltaK[Begriff des Monoidhomomorphismus]{%
  \MonoidHom{f}{A,\star,e}{B,\diamond,u}%
}{MonoidHomDef}{%
  \DeltaRow{Mengen}{A\dsep B\dsep e\dsep u}
  \DeltaRow{Funktionen}{f\dsep\star\dsep\diamond}
  \DeltaRow{\textbf{Axiome}}{}
  \DeltaPrem{Quellmonoid}{\MonoidAlg{A}{\star}{e}}
  \DeltaPrem{Zielmonoid}{\MonoidAlg{B}{\diamond}{u}}
  \DeltaPrem{Halbgruppenhomomorphismus}{%
    \SemigroupHom{f}{A,\star}{B,\diamond}}
  \DeltaRow{Erhaltung des neutralen Elements}{f(e)=u}
  \DeltaRow{\textbf{Neues Symbol}}{}
  \DeltaPrem{Monoidhomomorphismus}{%
    \MonoidHom{f}{A,\star,e}{B,\diamond,u}}
}

\FormulaAxiomDeltaK[Operationserhaltung eines Monoidhomomorphismus]{%
  \MonoidHom{f}{A,\star,e}{B,\diamond,u}\dsep x,y\in A
  \vdash f(x\star y)=f(x)\diamond f(y)%
}{MonoidHomOperationAxiom}{%
  \DeltaRow{Mengen}{A\dsep B\dsep e\dsep u\dsep x\dsep y}
  \DeltaRow{Funktionen}{f\dsep\star\dsep\diamond}
}

\FormulaAxiomDeltaK[Erhaltung des neutralen Elements]{%
  \MonoidHom{f}{A,\star,e}{B,\diamond,u}
  \vdash f(e)=u%
}{MonoidHomIdentityAxiom}{%
  \DeltaRow{Mengen}{A\dsep B\dsep e\dsep u}
  \DeltaRow{Funktionen}{f\dsep\star\dsep\diamond}
}

\FormulaDefDeltaK[Begriff des Monoidisomorphismus]{%
  \MonoidIso{f}{A,\star,e}{B,\diamond,u}%
}{MonoidIsoDef}{%
  \DeltaRow{Mengen}{A\dsep B\dsep e\dsep u}
  \DeltaRow{Funktionen}{f\dsep\star\dsep\diamond}
  \DeltaRow{\textbf{Axiome}}{}
  \DeltaPrem{Monoidhomomorphismus}{%
    \MonoidHom{f}{A,\star,e}{B,\diamond,u}}
  \DeltaPrem{Bijektivität}{f\colon A\bij B}
  \DeltaRow{\textbf{Neues Symbol}}{}
  \DeltaPrem{Monoidisomorphismus}{%
    \MonoidIso{f}{A,\star,e}{B,\diamond,u}}
}

\subsection{Transport durch Halbgruppenisomorphismen}

\ProofWideReasonsNearFormula

\FormulaThmDeltaK[Die Bilder neutraler Produkte]{%
  \begin{aligned}[t]
    &\MonoidAlg{A}{\star}{e}\dsep
      \SemigroupIso{F}{A,\star}{B,\diamond}\dsep a\in A\vdash{}\\[-2pt]
    &F(e)\diamond F(a)=F(a)\ \land\ F(a)\diamond F(e)=F(a)
  \end{aligned}
}{SemigroupIsoMonoidImageActions}{%
  \DeltaRow{Mengen}{A\dsep B\dsep e\dsep a}
  \DeltaRow{Funktionen}{F\dsep\star\dsep\diamond}
}
\begin{tabproofwide}
  \proofstepwidestarbreakkeep[1]{\MonoidAlg{A}{\star}{e}}{\rA}
  \proofstepwidestarbreakkeep[2]{\SemigroupIso{F}{A,\star}{B,\diamond}}{\rA}
  \proofstepwidestarbreakkeep[3]{a\in A}{\rA}
  \proofstepwidestarbreakkeep[1]{e\in A}{\FormulaRefAuto{MonoidIdentityCarrierAxiom}[ax]{1}}
  \proofstepwidestarbreakkeep[2]{\SemigroupHom{F}{A,\star}{B,\diamond}}{%
    \FormulaRefAuto{SemigroupIsoHomAxiom}[ax]{2}}
  \proofstepwidestarbreakkeep[1,3]{e\star a=a\dsep a\star e=a}{%
    \FormulaRefAuto{MonoidLeftIdentityAxiom}[ax]{1,3},\
    \FormulaRefAuto{MonoidRightIdentityAxiom}[ax]{1,3}}
  \proofstepwidestarbreakkeep[1,2,3]{F(e\star a)=F(e)\diamond F(a)}{%
    \FormulaRefAuto{SemigroupHomOperationAxiom}[ax]{5,4,3}}
  \proofstepwidestarbreakkeep[1,2,3]{F(a\star e)=F(a)\diamond F(e)}{%
    \FormulaRefAuto{SemigroupHomOperationAxiom}[ax]{5,3,4}}
  \proofstepwidestarbreakkeep[1,2,3]{F(e)\diamond F(a)=F(a)}{%
    \rIE{6,\FormulaRefAuto{a=b\vdash b=a}{7}}}
  \proofstepwidestarbreakkeep[1,2,3]{F(a)\diamond F(e)=F(a)}{%
    \rIE{6,\FormulaRefAuto{a=b\vdash b=a}{8}}}
  \proofstepwidestarbreakkeep[1,2,3]{F(e)\diamond F(a)=F(a)\ \land\ F(a)\diamond F(e)=F(a)}{\rAI{9,10}}
\end{tabproofwide}

\FormulaThmDeltaK[Ein Halbgruppenisomorphismus transportiert die Monoidstruktur]{%
  \MonoidAlg{A}{\star}{e}\dsep
  \SemigroupIso{F}{A,\star}{B,\diamond}
  \vdash\MonoidAlg{B}{\diamond}{F(e)}
}{SemigroupIsoMonoidTransport}{%
  \DeltaRow{Mengen}{A\dsep B\dsep e\dsep a\dsep b}
  \DeltaRow{Funktionen}{F\dsep\star\dsep\diamond}
}
\begin{tabproofsplitwide}
  \proofpartwide{Halbgruppenstruktur und Typisierung der Eins}
  \proofstepwidestarbreakkeep[1]{\MonoidAlg{A}{\star}{e}}{\rA}
  \proofstepwidestarbreakkeep[2]{\SemigroupIso{F}{A,\star}{B,\diamond}}{\rA}
  \proofstepwidestarbreakkeep[2]{\SemigroupHom{F}{A,\star}{B,\diamond}}{%
    \FormulaRefAuto{SemigroupIsoHomAxiom}[ax]{2}}
  \proofstepwidestarbreakkeep[2]{\SemigroupAlg{B}{\diamond}}{%
    \FormulaRefAuto{SemigroupHomTargetAxiom}[ax]{3}}
  \proofstepwidestarbreakkeep[2]{F\colon A\bij B}{%
    \FormulaRefAuto{SemigroupIsoBijectiveAxiom}[ax]{2}}
  \proofstepwidestarbreakkeep[1]{e\in A}{\FormulaRefAuto{MonoidIdentityCarrierAxiom}[ax]{1}}
  \proofstepwidestarbreakkeep[1,2]{F(e)\in B}{%
    \FormulaRefAuto{F\colon A\bij B\dsep x\in A\vdash F(x)\in B}{5,6}}
  \closeproofpartwide

  \proofpartwide{Beidseitige Neutralität im Ziel}
  \setcounter{proofstepnr}{7}
  \proofstepwidestarbreakkeep[8]{b\in B}{\rA}
  \proofstepwidestarbreakkeep[2,8]{\exists a\in A\,(F(a)=b)}{%
    \FormulaRefAuto{Bijektive Funktion}[def]{5,8}}
  \proofstepwidestarbreakkeep[10]{a\in A\land F(a)=b}{\rA}
  \proofstepwidestarbreakkeep[1,2,10]{F(e)\diamond F(a)=F(a)\land F(a)\diamond F(e)=F(a)}{%
    \FormulaRefAuto{SemigroupIsoMonoidImageActions}[thm]{1,2,\rAEa{10}}}
  \proofstepwidestarbreakkeep[1,2,10]{F(e)\diamond b=b\land b\diamond F(e)=b}{\rIE{\rAEb{10},11}}
  \proofstepwidestarbreakkeep[1,2,8]{F(e)\diamond b=b\land b\diamond F(e)=b}{\rEE{9,10,12}}
  \proofstepwidestarbreakkeep[1,2]{\MonoidAlg{B}{\diamond}{F(e)}}{%
    \FormulaRefAuto{MonoidDef}[def]{4,7,\rAEa{13},\rAEb{13}}}
  \closeproofpartwide
\end{tabproofsplitwide}

\FormulaThmDeltaK[Halbgruppenisomorphismen zwischen Monoiden erhalten die Eins]{%
  \MonoidAlg{A}{\star}{e}\dsep\MonoidAlg{B}{\diamond}{u}\dsep
  \SemigroupIso{F}{A,\star}{B,\diamond}\vdash F(e)=u
}{SemigroupIsoBetweenMonoidsIdentity}{%
  \DeltaRow{Mengen}{A\dsep B\dsep e\dsep u\dsep b}
  \DeltaRow{Funktionen}{F\dsep\star\dsep\diamond}
}
\begin{tabproofwide}
  \proofstepwidestarbreakkeep[1]{\MonoidAlg{A}{\star}{e}}{\rA}
  \proofstepwidestarbreakkeep[2]{\MonoidAlg{B}{\diamond}{u}}{\rA}
  \proofstepwidestarbreakkeep[3]{\SemigroupIso{F}{A,\star}{B,\diamond}}{\rA}
  \proofstepwidestarbreakkeep[1,3]{\MonoidAlg{B}{\diamond}{F(e)}}{%
    \FormulaRefAuto{SemigroupIsoMonoidTransport}[thm]{1,3}}
  \proofstepwidestarbreakkeep[2]{u\in B}{\FormulaRefAuto{MonoidIdentityCarrierAxiom}[ax]{2}}
  \proofstepwidestarbreakkeep[6]{b\in B}{\rA}
  \proofstepwidestarbreakkeep[2,6]{u\diamond b=b}{\FormulaRefAuto{MonoidLeftIdentityAxiom}[ax]{2,6}}
  \proofstepwidestarbreakkeep[2,6]{b\diamond u=b}{\FormulaRefAuto{MonoidRightIdentityAxiom}[ax]{2,6}}
  \proofstepwidestarbreakkeep[1,2,3]{F(e)=u}{%
    \FormulaRefAuto{MonoidIdentityCandidateEqualsIdentity}[thm]{4,5,\rUI{\rRI{6,7}},\rUI{\rRI{6,8}}}}
\end{tabproofwide}

\FormulaThmDeltaK[Ein Halbgruppenisomorphismus erhält und reflektiert inverse Paare]{%
  \begin{aligned}[t]
    &\MonoidAlg{A}{\star}{e}\dsep\MonoidAlg{B}{\diamond}{u}\dsep
      \SemigroupIso{F}{A,\star}{B,\diamond}\dsep a,b\in A\vdash{}\\[-2pt]
    &(a\star b=e\land b\star a=e)\leftrightarrow
      (F(a)\diamond F(b)=u\land F(b)\diamond F(a)=u)
  \end{aligned}
}{SemigroupIsoMonoidInversePair}{%
  \DeltaRow{Mengen}{A\dsep B\dsep e\dsep u\dsep a\dsep b}
  \DeltaRow{Funktionen}{F\dsep\star\dsep\diamond}
}
\begin{tabproofsplitwide}
  \proofpartwide{Transportdaten und Erhaltung inverser Paare}
  \proofstepwidestarbreakkeep[1]{\MonoidAlg{A}{\star}{e}}{\rA}
  \proofstepwidestarbreakkeep[2]{\MonoidAlg{B}{\diamond}{u}}{\rA}
  \proofstepwidestarbreakkeep[3]{\SemigroupIso{F}{A,\star}{B,\diamond}}{\rA}
  \proofstepwidestarbreakkeep[4]{a,b\in A}{\rA}
  \proofstepwidestarbreakkeep[3]{\SemigroupHom{F}{A,\star}{B,\diamond}}{%
    \FormulaRefAuto{SemigroupIsoHomAxiom}[ax]{3}}
  \proofstepwidestarbreakkeep[1,2,3]{F(e)=u}{%
    \FormulaRefAuto{SemigroupIsoBetweenMonoidsIdentity}[thm]{1,2,3}}
  \proofstepwidestarbreakkeep[3,4]{F(a\star b)=F(a)\diamond F(b)}{%
    \FormulaRefAuto{SemigroupHomOperationAxiom}[ax]{5,4}}
  \proofstepwidestarbreakkeep[3,4]{F(b\star a)=F(b)\diamond F(a)}{%
    \FormulaRefAuto{SemigroupHomOperationAxiom}[ax]{5,4}}
  \proofstepwidestar[3]{F\colon A\bij B}{%
      \FormulaRefAuto{SemigroupIsoBijectiveAxiom}[ax]{3}}
  \proofstepwidestarbreakkeep[3]{F\colon A\inj B}{%
    \FormulaRefAuto{F\colon A\bij B\vdash F\colon A\inj B}{%
      9}}
  \proofstepwidestar[1,4]{\SemigroupAlg{A}{\star}}{%
      \FormulaRefAuto{MonoidBaseSemigroupAxiom}[ax]{1}}
  \proofstepwidestarbreakkeep[1,4]{a\star b,b\star a\in A}{%
    \FormulaRefAuto{SemigroupClosureAxiom}[ax]{%
      11,4}}
  \proofstepwidestarbreakkeep[1]{e\in A}{\FormulaRefAuto{MonoidIdentityCarrierAxiom}[ax]{1}}
  \proofstepwidestarbreakkeep[14]{a\star b=e\land b\star a=e}{\rA}
  \proofstepwidestarbreakkeep[1,2,3,4,14]{F(a)\diamond F(b)=u\land F(b)\diamond F(a)=u}{%
    \rAI{\rIE{\rAEa{14},6,7},\rIE{\rAEb{14},6,8}}}
  \closeproofpartwide

  \proofpartwide{Reflexion inverser Paare}
  \setcounter{proofstepnr}{15}
  \proofstepwidestarbreakkeep[16]{F(a)\diamond F(b)=u\land F(b)\diamond F(a)=u}{\rA}
  \proofstepwidestarbreakkeep[1,2,3,4,16]{F(a\star b)=F(e)\dsep F(b\star a)=F(e)}{%
    \rIE{6,7,8,16}}
  \proofstepwidestarbreakkeep[1,2,3,4,16]{a\star b=e}{%
    \FormulaRefAuto{F\colon A\inj B\dsep x\in A\dsep y\in A\dsep F(x)=F(y)\vdash x=y}[ax]{10,12,13,17}}
  \proofstepwidestarbreakkeep[1,2,3,4,16]{b\star a=e}{%
    \FormulaRefAuto{F\colon A\inj B\dsep x\in A\dsep y\in A\dsep F(x)=F(y)\vdash x=y}[ax]{10,12,13,17}}
  \proofstepwidestarbreakkeep[1,2,3,4]{%
    (a\star b=e\land b\star a=e)\leftrightarrow
    (F(a)\diamond F(b)=u\land F(b)\diamond F(a)=u)}{%
    \rLRI{\rRI{14,15},\rRI{16,\rAI{18,19}}}}
  \closeproofpartwide
\end{tabproofsplitwide}

\par\noindent\begin{minipage}{\linewidth}
\FormulaThmDeltaK[Halbgruppenisomorphismen zwischen Monoiden erhalten Einheiten]{%
  \begin{aligned}[t]
    &\MonoidAlg{A}{\star}{e}\dsep\MonoidAlg{B}{\diamond}{u}\dsep
      \SemigroupIso{F}{A,\star}{B,\diamond}\dsep a\in A\vdash{}\\[-2pt]
    &\MonoidUnit{A,\star,e}{a}\leftrightarrow
      \MonoidUnit{B,\diamond,u}{F(a)}
  \end{aligned}
}{SemigroupIsoBetweenMonoidsUnits}{%
  \DeltaRow{Mengen}{A\dsep B\dsep e\dsep u\dsep a\dsep b\dsep v}
  \DeltaRow{Funktionen}{F\dsep\star\dsep\diamond}
}
\end{minipage}\par
\begin{tabproofsplitwide}
  \proofpartwide{Erhaltung der Einheiten}
  \proofstepwidestarbreakkeep[1]{\MonoidAlg{A}{\star}{e}}{\rA}
  \proofstepwidestarbreakkeep[2]{\MonoidAlg{B}{\diamond}{u}}{\rA}
  \proofstepwidestarbreakkeep[3]{\SemigroupIso{F}{A,\star}{B,\diamond}}{\rA}
  \proofstepwidestarbreakkeep[4]{a\in A}{\rA}
  \proofstepwidestarbreakkeep[3]{F\colon A\bij B}{%
    \FormulaRefAuto{SemigroupIsoBijectiveAxiom}[ax]{3}}
  \proofstepwidestarbreakkeep[3,4]{F(a)\in B}{%
    \FormulaRefAuto{F\colon A\bij B\dsep x\in A\vdash F(x)\in B}{5,4}}
  \proofstepwidestarbreakkeep[7]{\MonoidUnit{A,\star,e}{a}}{\rA}
  \proofstepwidestarbreakkeep[7]{\exists b\in A\,(a\star b=e\land b\star a=e)}{%
    \FormulaRefAuto{MonoidUnitDef}[def]{7}}
  \proofstepwidestarbreakkeep[9]{b\in A\land(a\star b=e\land b\star a=e)}{\rA}
  \proofstepwidestarbreakkeep[3,9]{F(b)\in B}{%
    \FormulaRefAuto{F\colon A\bij B\dsep x\in A\vdash F(x)\in B}{5,\rAEa{9}}}
  \proofstepwidestarbreakkeep[1,2,3,4,9]{F(a)\diamond F(b)=u\land F(b)\diamond F(a)=u}{%
    \FormulaRefAuto{SemigroupIsoMonoidInversePair}[thm]{1,2,3,4,\rAEa{9},\rAEb{9}}}
  \proofstepwidestarbreakkeep[1,2,3,4,9]{\MonoidUnit{B,\diamond,u}{F(a)}}{%
    \FormulaRefAuto{MonoidUnitDef}[def]{2,6,\rEI{10,11}}}
  \proofstepwidestarbreakkeep[1,2,3,4,7]{\MonoidUnit{B,\diamond,u}{F(a)}}{\rEE{8,9,12}}
  \closeproofpartwide

  \proofpartwide{Reflexion der Einheiten}
  \setcounter{proofstepnr}{13}
  \proofstepwidestarbreakkeep[14]{\MonoidUnit{B,\diamond,u}{F(a)}}{\rA}
  \proofstepwidestarbreakkeep[14]{\exists v\in B\,(F(a)\diamond v=u\land v\diamond F(a)=u)}{%
    \FormulaRefAuto{MonoidUnitDef}[def]{14}}
  \proofstepwidestarbreakkeep[16]{v\in B\land(F(a)\diamond v=u\land v\diamond F(a)=u)}{\rA}
  \proofstepwidestarbreakkeep[3,16]{\exists b\in A\,(F(b)=v)}{%
    \FormulaRefAuto{Bijektive Funktion}[def]{5,\rAEa{16}}}
  \proofstepwidestarbreakkeep[18]{b\in A\land F(b)=v}{\rA}
  \proofstepwidestarbreakkeep[16,18]{F(a)\diamond F(b)=u\land F(b)\diamond F(a)=u}{%
    \rIE{\rAEb{18},\rAEb{16}}}
  \proofstepwidestarbreakkeep[1,2,3,4,16,18]{a\star b=e\land b\star a=e}{%
    \FormulaRefAuto{SemigroupIsoMonoidInversePair}[thm]{1,2,3,4,\rAEa{18},19}}
  \proofstepwidestarbreakkeep[1,2,3,4,16,18]{\MonoidUnit{A,\star,e}{a}}{%
    \FormulaRefAuto{MonoidUnitDef}[def]{1,4,\rEI{\rAEa{18},20}}}
  \proofstepwidestarbreakkeep[1,2,3,4,16]{\MonoidUnit{A,\star,e}{a}}{\rEE{17,18,21}}
  \proofstepwidestarbreakkeep[1,2,3,4,14]{\MonoidUnit{A,\star,e}{a}}{\rEE{15,16,22}}
  \proofstepwidestarbreakkeep[1,2,3,4]{\MonoidUnit{A,\star,e}{a}\leftrightarrow\MonoidUnit{B,\diamond,u}{F(a)}}{%
    \rLRI{\rRI{7,13},\rRI{14,23}}}
  \closeproofpartwide
\end{tabproofsplitwide}

\par\noindent\begin{minipage}{\linewidth}
\FormulaThmDeltaK[Identität und Einheiten unter einem Halbgruppenisomorphismus]{%
  \begin{aligned}[t]
    &\MonoidAlg{A}{\star}{e}\dsep\MonoidAlg{B}{\diamond}{u}\dsep
      \SemigroupIso{F}{A,\star}{B,\diamond}\vdash{}\\[-2pt]
    &F(e)=u\ \land\ \forall a\in A\,
      \bigl(\MonoidUnit{A,\star,e}{a}\leftrightarrow
        \MonoidUnit{B,\diamond,u}{F(a)}\bigr)
  \end{aligned}
}{SemigroupIsoBetweenMonoidsIdentityAndUnits}{%
  \DeltaRow{Mengen}{A\dsep B\dsep e\dsep u\dsep a}
  \DeltaRow{Funktionen}{F\dsep\star\dsep\diamond}
}
\end{minipage}\par
\begin{tabproofwide}
  \proofstepwidestarbreakkeep[1]{\MonoidAlg{A}{\star}{e}}{\rA}
  \proofstepwidestarbreakkeep[2]{\MonoidAlg{B}{\diamond}{u}}{\rA}
  \proofstepwidestarbreakkeep[3]{\SemigroupIso{F}{A,\star}{B,\diamond}}{\rA}
  \proofstepwidestarbreakkeep[1,2,3]{F(e)=u}{%
    \FormulaRefAuto{SemigroupIsoBetweenMonoidsIdentity}[thm]{1,2,3}}
  \proofstepwidestarbreakkeep[5]{a\in A}{\rA}
  \proofstepwidestarbreakkeep[1,2,3,5]{\MonoidUnit{A,\star,e}{a}\leftrightarrow\MonoidUnit{B,\diamond,u}{F(a)}}{%
    \FormulaRefAuto{SemigroupIsoBetweenMonoidsUnits}[thm]{1,2,3,5}}
  \proofstepwidestarbreakkeep[1,2,3]{F(e)=u\ \land\ \forall a\in A\,
    (\MonoidUnit{A,\star,e}{a}\leftrightarrow\MonoidUnit{B,\diamond,u}{F(a)})}{%
    \rAI{4,\rUI{\rRI{5,6}}}}
\end{tabproofwide}

\ProofWideReasonsAtRight

\subsection{Grundlegende Eigenschaften}

\FormulaThmDeltaK[Komposition von Monoidhomomorphismen]{%
  \begin{aligned}[t]
    &\MonoidHom{f}{A,\star,e}{B,\diamond,u}\dsep
      \MonoidHom{g}{B,\diamond,u}{C,\bullet,v}\\
    &\vdash\MonoidHom{g\circ f}{A,\star,e}{C,\bullet,v}
  \end{aligned}
}{MonoidHomComposition}{%
  \DeltaRow{Mengen}{A\dsep B\dsep C\dsep e\dsep u\dsep v}
  \DeltaRow{Funktionen}{f\dsep g\dsep\star\dsep\diamond\dsep\bullet}
}
\begin{tabproofsplitwide}
  \proofpartwide{Erhaltung der Halbgruppenoperation}
  \proofstepwidestar[1]{\MonoidHom{f}{A,\star,e}{B,\diamond,u}}{\rA}
  \proofstepwidestar[2]{\MonoidHom{g}{B,\diamond,u}{C,\bullet,v}}{\rA}
  \proofstepwidestar[1]{\SemigroupHom{f}{A,\star}{B,\diamond}}{%
    \FormulaRefAuto{MonoidHomDef}[def]{1}}
  \proofstepwidestar[2]{\SemigroupHom{g}{B,\diamond}{C,\bullet}}{%
    \FormulaRefAuto{MonoidHomDef}[def]{2}}
  \proofstepwidestar[1,2]{%
    \SemigroupHom{g\circ f}{A,\star}{C,\bullet}}{%
    \FormulaRefAuto{SemigroupHomComposition}[thm]{3,4}}
  \closeproofpartwide

  \proofpartwide{Erhaltung der Eins und Monoidstruktur}
  \setcounter{proofstepnr}{5}
  \proofstepwidestar[1]{f(e)=u}{%
    \FormulaRefAuto{MonoidHomIdentityAxiom}[ax]{1}}
  \proofstepwidestar[2]{g(u)=v}{%
    \FormulaRefAuto{MonoidHomIdentityAxiom}[ax]{2}}
  \proofstepwidestar[1]{\MonoidAlg{A}{\star}{e}}{%
    \FormulaRefAuto{MonoidHomDef}[def]{1}}
  \proofstepwidestar[1]{e\in A}{%
    \FormulaRefAuto{MonoidIdentityCarrierAxiom}[ax]{8}}
  \proofstepwidestar[1]{f\colon A\to B}{%
    \FormulaRefAuto{SemigroupHomFunctionAxiom}[ax]{3}}
  \proofstepwidestar[2]{g\colon B\to C}{%
    \FormulaRefAuto{SemigroupHomFunctionAxiom}[ax]{4}}
  \proofstepwide[1,2]{(g\circ f)(e)}{=}{g(f(e))}{%
    \FormulaRefAuto{x\in A \vdash (G\circ F)(x)=G(F(x))}{%
      10,11,9}}
  \proofstepwide[1,2]{}{=}{g(u)}{\rIE{6}}
  \proofstepwide[1,2]{}{=}{v}{\rIE{7}}
  \proofstepwidestar[1,2]{(g\circ f)(e)=v}{\rChain{12,14}}
  \proofstepwidestar[2]{\MonoidAlg{C}{\bullet}{v}}{%
    \FormulaRefAuto{MonoidHomDef}[def]{2}}
  \proofstepwidestar[1,2]{%
    \MonoidHom{g\circ f}{A,\star,e}{C,\bullet,v}}{%
    \FormulaRefAuto{MonoidHomDef}[def]{8,16,5,15}}
  \closeproofpartwide
\end{tabproofsplitwide}

\subsection{Kommutative Homomorphismen und Isomorphismen}

\FormulaDefDeltaK[Homomorphismen kommutativer Monoide]{%
  \CommutativeMonoidHom{f}{A,\star,e}{B,\diamond,u}%
}{CommutativeMonoidHomDef}{%
  \DeltaRow{Mengen}{A\dsep B\dsep e\dsep u}
  \DeltaRow{Funktionen}{f\dsep\star\dsep\diamond}
  \DeltaRow{\textbf{Axiome}}{}
  \DeltaPrem{Quellstruktur}{\CommutativeMonoidAlg{A}{\star}{e}}
  \DeltaPrem{Zielstruktur}{\CommutativeMonoidAlg{B}{\diamond}{u}}
  \DeltaPrem{Homomorphismus}{\MonoidHom{f}{A,\star,e}{B,\diamond,u}}
  \DeltaRow{\textbf{Neues Symbol}}{}
  \DeltaPrem{Homomorphismus}{\CommutativeMonoidHom{f}{A,\star,e}{B,\diamond,u}}
}

\FormulaAxiomDeltaK[Quellstruktur eines Homomorphismus kommutativer Monoide]{%
  \CommutativeMonoidHom{f}{A,\star,e}{B,\diamond,u}
  \vdash\CommutativeMonoidAlg{A}{\star}{e}%
}{CommutativeMonoidHomSourceAxiom}{%
  \DeltaRow{Mengen}{A\dsep B\dsep e\dsep u}
  \DeltaRow{Funktionen}{f\dsep\star\dsep\diamond}
}

\FormulaAxiomDeltaK[Zielstruktur eines Homomorphismus kommutativer Monoide]{%
  \CommutativeMonoidHom{f}{A,\star,e}{B,\diamond,u}
  \vdash\CommutativeMonoidAlg{B}{\diamond}{u}%
}{CommutativeMonoidHomTargetAxiom}{%
  \DeltaRow{Mengen}{A\dsep B\dsep e\dsep u}
  \DeltaRow{Funktionen}{f\dsep\star\dsep\diamond}
}

\FormulaAxiomDeltaK[Monoidhomomorphismus eines kommutativen Homomorphismus]{%
  \CommutativeMonoidHom{f}{A,\star,e}{B,\diamond,u}
  \vdash\MonoidHom{f}{A,\star,e}{B,\diamond,u}%
}{CommutativeMonoidHomBaseHomAxiom}{%
  \DeltaRow{Mengen}{A\dsep B\dsep e\dsep u}
  \DeltaRow{Funktionen}{f\dsep\star\dsep\diamond}
}

\FormulaDefDeltaK[Isomorphismen kommutativer Monoide]{%
  \CommutativeMonoidIso{f}{A,\star,e}{B,\diamond,u}%
}{CommutativeMonoidIsoDef}{%
  \DeltaRow{Mengen}{A\dsep B\dsep e\dsep u}
  \DeltaRow{Funktionen}{f\dsep\star\dsep\diamond}
  \DeltaRow{\textbf{Axiome}}{}
  \DeltaPrem{Homomorphismus}{\CommutativeMonoidHom{f}{A,\star,e}{B,\diamond,u}}
  \DeltaPrem{Bijektivität}{f\colon A\bij B}
  \DeltaRow{\textbf{Neues Symbol}}{}
  \DeltaPrem{Isomorphismus}{\CommutativeMonoidIso{f}{A,\star,e}{B,\diamond,u}}
}

\FormulaAxiomDeltaK[Homomorphismus eines Isomorphismus kommutativer Monoide]{%
  \CommutativeMonoidIso{f}{A,\star,e}{B,\diamond,u}
  \vdash\CommutativeMonoidHom{f}{A,\star,e}{B,\diamond,u}%
}{CommutativeMonoidIsoHomAxiom}{%
  \DeltaRow{Mengen}{A\dsep B\dsep e\dsep u}
  \DeltaRow{Funktionen}{f\dsep\star\dsep\diamond}
}

\FormulaAxiomDeltaK[Bijektivität eines Isomorphismus kommutativer Monoide]{%
  \CommutativeMonoidIso{f}{A,\star,e}{B,\diamond,u}
  \vdash f\colon A\bij B%
}{CommutativeMonoidIsoBijectiveAxiom}{%
  \DeltaRow{Mengen}{A\dsep B\dsep e\dsep u}
  \DeltaRow{Funktionen}{f\dsep\star\dsep\diamond}
}

\chapter{Beispiele von Monoiden}

\section{Die natürlichen Zahlen}

\FormulaThmDeltaK[Die Addition bildet ein Monoid]{%
  \MonoidAlg{\mathbb{N}}{+}{0}%
}{NaturalAdditionMonoid}{%
  \DeltaRow{Mengen}{\mathbb{N}\dsep 0\dsep a}
}
\begin{tabproof}
  \proofstep{}{\SemigroupAlg{\mathbb{N}}{+}}{%
    \FormulaRefAuto{NaturalAdditionSemigroup}[thm]}
  \proofstep{}{0\in\mathbb{N}}{%
    \FormulaRefAuto{PeanoZero}[ax]}
  \proofstep{3}{a\in\mathbb{N}}{\rA}
  \proofstep{3}{0+a=a}{%
    \FormulaRefAuto{PeanoAddLeftZero}[thm]{3}}
  \proofstep{3}{a+0=a}{%
    \FormulaRefAuto{PeanoAddRightZero}[thm]{3}}
  \proofstep{}{\MonoidAlg{\mathbb{N}}{+}{0}}{%
    \FormulaRefAuto{MonoidDef}[def]{1,2,4,5}}
\end{tabproof}

\FormulaThmDeltaK[Die Addition bildet ein kommutatives Monoid]{%
  \CommutativeMonoidAlg{\mathbb{N}}{+}{0}%
}{NaturalAdditiveCommutativeMonoid}{%
  \DeltaRow{Mengen}{\mathbb{N}\dsep 0\dsep a\dsep b}
}
\begin{tabproof}
  \proofstep{}{\MonoidAlg{\mathbb{N}}{+}{0}}{%
    \FormulaRefAuto{NaturalAdditionMonoid}[thm]}
  \proofstep{2}{a\in\mathbb{N}}{\rA}
  \proofstep{3}{b\in\mathbb{N}}{\rA}
  \proofstep{2,3}{a+b=b+a}{%
    \FormulaRefAuto{PeanoAddCommutative}[thm]{2,3}}
  \proofstep{}{\CommutativeMonoidAlg{\mathbb{N}}{+}{0}}{%
    \FormulaRefAuto{CommutativeMonoidDef}[def]{1,4}}
\end{tabproof}

\FormulaThmDeltaK[Die Multiplikation bildet ein Monoid]{%
  \MonoidAlg{\mathbb{N}}{\cdot}{1}%
}{NaturalMultiplicationMonoid}{%
  \DeltaRow{Mengen}{\mathbb{N}\dsep 1\dsep a}
}
\begin{tabproof}
  \proofstep{}{\SemigroupAlg{\mathbb{N}}{\cdot}}{%
    \FormulaRefAuto{NaturalMultiplicationSemigroup}[thm]}
  \proofstep{}{1\in\mathbb{N}}{%
    \FormulaRefAuto{PeanoOneInN}[thm]}
  \proofstep{3}{a\in\mathbb{N}}{\rA}
  \proofstep{3}{1\cdot a=a}{%
    \FormulaRefAuto{PeanoMulLeftOne}[thm]{3}}
  \proofstep{3}{a\cdot1=a}{%
    \FormulaRefAuto{PeanoMulRightOne}[thm]{3}}
  \proofstep{}{\MonoidAlg{\mathbb{N}}{\cdot}{1}}{%
    \FormulaRefAuto{MonoidDef}[def]{1,2,4,5}}
\end{tabproof}

\FormulaThmDeltaK[Die Multiplikation bildet ein kommutatives Monoid]{%
  \CommutativeMonoidAlg{\mathbb{N}}{\cdot}{1}%
}{NaturalMultiplicativeCommutativeMonoid}{%
  \DeltaRow{Mengen}{\mathbb{N}\dsep 1\dsep a\dsep b}
}
\begin{tabproof}
  \proofstep{}{\MonoidAlg{\mathbb{N}}{\cdot}{1}}{%
    \FormulaRefAuto{NaturalMultiplicationMonoid}[thm]}
  \proofstep{2}{a\in\mathbb{N}}{\rA}
  \proofstep{3}{b\in\mathbb{N}}{\rA}
  \proofstep{2,3}{a\cdot b=b\cdot a}{%
    \FormulaRefAuto{PeanoMulCommutative}[thm]{2,3}}
  \proofstep{}{\CommutativeMonoidAlg{\mathbb{N}}{\cdot}{1}}{%
    \FormulaRefAuto{CommutativeMonoidDef}[def]{1,4}}
\end{tabproof}

\end{document}
